\documentclass[UTF8, 12pt, fontset=fandol, oneside]{ctexart}
\usepackage{researchnotes}

% 保留分册独立编译时的章号前缀。
\renewcommand{\thesection}{2.\arabic{section}}

\title{第二章\quad 序列的极限}
\author{}
\date{}

\begin{document}

\maketitle

极限是构筑微积分坚实理论体系的基石. 要想对数学分析这门学科的实质有一个真正的了解和掌握, 就必须准确掌握极限的概念和无穷小的分析方法. 这一章我们将致力于序列极限的讨论, 从中建立极限的主要理论与方法. 下一章我们来讨论函数的极限。

\section{序列极限的定义}

\subsection{序列}

所谓的\textbf{序列} (有时也称为\textbf{数列}) 实质上就是一个从正整数集 $\mathbb{N}$ 到实数集 $\mathbb{R}$ 的一个函数 $f:\mathbb{N}\to\mathbb{R}$. 但我们常将一个序列看做是按照一定顺序排列的一列数:
\[
  x_1=f(1),\ x_2=f(2),\ \cdots,\ x_n=f(n),\ \cdots .
\]
一个序列通常记为 $\{x_n\}$, 其中 $x_n$ 称为\textbf{通项}. 有时为了简单起见, 我们也将构成一个序列的所有数作成的集合记为 $\{x_n\}$, 即此时 $\{x_n\}=f(\mathbb{N})$. 读者以后应注意``集合 $\{x_n\}$''和``序列 $\{x_n\}$''的区别。

例如, 序列
\[
  1,\ \frac12,\ \frac13,\ \cdots,\ \frac1n,\ \cdots
\]
的通项是 $\dfrac1n$. 但在有的时候序列的通项却很难用公式给出, 如圆周率 $\pi$ 精确到小数点后 $n$ 位的有限小数组成的序列
\[
  3,\ 3.1,\ 3.14,\ 3.141,\ 3.1415,\ \cdots
\]
就不可能写出通项公式。

我们已经知道在任意两个实数之间必有有理数, 有理数集的这一性质称为有理数集在实数集中的稠密性. 使人吃惊的是, 在数轴上密密麻麻分布的有理数的全体竟然可以排成一个序列. 下面我们仅排列 $[0,1)$ 中的有理数, $\mathbb{R}$ 中有理数的排列留给读者作为练习。

由于 $(0,1)$ 中的任何一个有理数都可唯一地表示为 $\dfrac pq$, 其中 $p$ 和 $q$ 是两个互素的正整数且 $p<q$, 因此我们可以将 $[0,1)$ 中的所有有理数按照其分母的大小从小到大排列, 而相同分母的再按其分子的大小从小到大排列. 这样, 我们就可得到序列:
\[
  0,\ \frac12,\ \frac13,\ \frac23,\ \frac14,\ \frac34,\ \frac15,\ \frac25,\ \frac35,\ \cdots .
\]
对此序列, 似乎也很难写出它的通项公式。

\subsection{序列极限的定义}

对于给定的一个序列 $\{x_n\}$, 我们这里主要关心的是随着 $n$ 的增大, $x_n$ 的变化趋势. 为此, 我们首先必须说清楚``变化趋势''是什么意思. 下面我们先看几个具体的例子。

\noindent\textbf{例 2.1.1}\quad 序列 $x_n=\dfrac1n(n=1,2,\cdots)$ 随着 $n$ 的增大, $x_n$ 从 $0$ 的右边越来越趋近于 $0$。

\noindent\textbf{例 2.1.2}\quad 序列 $x_n=\dfrac{(-1)^n}{n}(n=1,2,\cdots)$ 随着 $n$ 的增大, $x_n$ 时而在 $0$ 点的右边, 时而在 $0$ 点的左边, 但 $x_n$ 与 $0$ 的距离越来越趋近于 $0$。

\noindent\textbf{例 2.1.3}\quad 序列 $\{x_n\}$ 由下列法则确定:
\[
  x_{2n}=\frac1{2n},\qquad
  x_{2n+1}=\frac1{2^{2n+1}}
  \quad(n=1,2,\cdots).
\]
对于这一序列, 虽然我们不能说, 随着 $n$ 的增大, $x_n$ 与 $0$ 的距离越来越趋近于 $0$, 但 $x_n$ 与 $0$ 的距离也是无限地接近于 $0$, 只不过奇数项与偶数项接近零的速度不同而已。

\noindent\textbf{例 2.1.4}\quad 设 $\{x_n\}$ 由下式给出:
\[
  x_{2n-1}=\frac1n,\qquad
  x_{2n}=0
  \quad(n=1,2,\cdots),
\]
则 $\{x_n\}$ 有无穷多项为 $0$, 而随着 $n$ 的增大, $x_{2n-1}$ 也无限地趋于 $0$. 因此, 随着 $n$ 的增大, 整个序列也是趋近于 $0$ 的。

分析以上例子我们发现, 当 $n$ 无限变大时, 尽管它们的变化过程各具特点, 但四个序列最终都无限地趋于一个常数 $a=0$. 这种随着 $n$ 的增大可以无限地趋于一个常数的序列正是我们感兴趣的. 那么怎样来精确地描述一个序列 $\{x_n\}$ 无限地趋于一个常数 $a$ 呢? 也许大家会说, 这是指 $x_n$ 与 $a$ 的差 (随着 $n$ 的增大) 可以达到任意小的程度. 显然, 这只是一种描述性的说法, 十分不宜于进行严格的数学推理和演绎. 其实, 给出某精确的定义并非一件易事, 经过众多数学家的不懈努力和不断探索, 直到 19 世纪才有了数学上的如下定义:

\noindent\textbf{定义 2.1.1}\quad 设 $\{x_n\}$ 是一个序列. 若存在常数 $a\in\mathbb{R}$, 使得 $\forall\varepsilon>0$, $\exists N\in\mathbb{N}$, 当 $n>N$ 时, 有
\[
  |x_n-a|<\varepsilon,
\]
则称该序列是\textbf{收敛的}, 并称 $a$ 为该序列的\textbf{极限} (或者说序列 $\{x_n\}$ 收敛于 $a$), 记做 $\lim\limits_{n\to\infty}x_n=a$ 或 $x_n\to a(n\to\infty)$. 若不存在 $a\in\mathbb{R}$, 使得 $\{x_n\}$ 收敛于 $a$, 则称之为\textbf{发散序列}。

在上述定义中, $\varepsilon$ 是事先给定的任意小的正数, 若对 $\varepsilon_1>0$, 我们已经找到了正整数 $N_1$, 使得当 $n>N_1$ 时, 有 $|x_n-a|<\varepsilon_1$, 则对 $\forall\varepsilon_2>\varepsilon_1$, 当 $n>N_1$ 时, 必有 $|x_n-a|<\varepsilon_1<\varepsilon_2$. 所以, $\varepsilon$ 贵在``小''. 此外, 一般来说, $N$ 是依赖于 $\varepsilon$ 的, $\varepsilon$ 越小, 所需的 $N$ 就会越大. 对一个给定的正数 $\varepsilon$, 如果 $N$ 是满足定义中的要求的正整数, 则显然对任何比 $N$ 大的正整数也满足定义的要求。

现在我们来考查序列极限的几何意义. 极限定义中的不等式
\[
  |x_n-a|<\varepsilon,\qquad \forall n>N,
\]
可以写成
\[
  a-\varepsilon<x_n<a+\varepsilon,\qquad \forall n>N,
\]
即
\[
  x_n\in U(a,\varepsilon)=(a-\varepsilon,a+\varepsilon),\qquad \forall n>N.
\]
因此, 如果采用几何的语言, 极限的定义可以表述为: $\forall\varepsilon>0$, 在 $a$ 的 $\varepsilon$ 邻域 $U(a,\varepsilon)$ 内包含了 $\{x_n\}$ 自某项之后的所有项. 与之等价说法是: $\forall\varepsilon>0$, 在 $a$ 的 $\varepsilon$ 邻域 $U(a,\varepsilon)$ 外只有 $\{x_n\}$ 的有限项. 如图 2.1.1 所示。

\noindent\textbf{思考题}\quad 将序列 $\{x_n\}$ 看成是定义在正整数集上的函数, 试描述序列极限的几何意义。

\noindent\textbf{例 2.1.5}\quad 证明 $\lim\limits_{n\to\infty}\dfrac{n}{n+1}=1$。

\noindent\textbf{证明}\quad 对任意给定的 $\varepsilon>0$, 要使
\[
  \left|\frac{n}{n+1}-1\right|
  =
  \frac1{n+1}<\varepsilon,
\]
只需
\[
  n>\frac1\varepsilon-1.
\]
于是, 取大于 $\dfrac1\varepsilon-1$ 的任意正整数作为 $N$ (例如, 取 $N=\left[\dfrac1\varepsilon\right]$), 则当 $n>N$ 时, 就有
\[
  \left|\frac{n}{n+1}-1\right|
  =
  \frac1{n+1}<\varepsilon.
\]
由极限定义知, $\lim\limits_{n\to\infty}\dfrac{n}{n+1}=1$ 成立。

\noindent\textbf{例 2.1.6}\quad 证明 $\lim\limits_{n\to\infty}q^n=0\ (|q|<1)$。

\noindent\textbf{证明}\quad 不妨设 $q\ne0$, 否则该序列为常序列 $(x_n=0,\ n=1,2,\cdots)$, 所述极限自然成立。

$\forall\varepsilon>0$ (不妨设 $\varepsilon<1$), 要使
\[
  |q^n-0|=|q|^n<\varepsilon,
\]
只要
\[
  n\ln |q|<\ln\varepsilon.
\]
注意, 这里 $\ln |q|<0,\ \ln\varepsilon<0$, 于是上式等价于
\[
  n>\frac{\ln\varepsilon}{\ln |q|}.
\]
因此, 取 $N=\left[\dfrac{\ln\varepsilon}{\ln |q|}\right]$, 则当 $n>N$ 时, 就有 $|q^n-0|<\varepsilon$, 即 $\lim\limits_{n\to\infty}q^n=0$ 成立。

\noindent\textbf{例 2.1.7}\quad 证明 $\lim\limits_{n\to\infty}\dfrac{n^2-n+2}{3n^2+2n+4}=\dfrac13$。

\noindent\textbf{证明}\quad 对任意的正整数 $n$, 我们有
\[
  \left|\frac{n^2-n+2}{3n^2+2n+4}-\frac13\right|
  =
  \frac{5n-2}{3(3n^2+2n+4)}
  <
  \frac{5n}{9n^2}
  <
  \frac1n.
\]
于是, $\forall\varepsilon>0$, 只要取 $N=\left[\dfrac1\varepsilon\right]+1$, 则当 $n>N$ 时, 就有
\[
  \left|\frac{n^2-n+2}{3n^2+2n+4}-\frac13\right|
  <
  \frac1n<\varepsilon.
\]
这就证明了 $\lim\limits_{n\to\infty}\dfrac{n^2-n+2}{3n^2+2n+4}=\dfrac13$。

\noindent\textbf{例 2.1.8}\quad 证明 $\lim\limits_{n\to\infty}\sqrt[n]{a}=1(a>1)$。

\noindent\textbf{证法 1}\quad $\forall\varepsilon>0$, 要使
\[
  |\sqrt[n]{a}-1|=\sqrt[n]{a}-1<\varepsilon,
\]
只要 $\sqrt[n]{a}<1+\varepsilon$, 即 $\dfrac1n\lg a<\lg(1+\varepsilon)$, 这只要 $n>\dfrac{\lg a}{\lg(1+\varepsilon)}$. 于是, 取 $N=\left[\dfrac{\lg a}{\lg(1+\varepsilon)}\right]$, 则当 $n>N$ 时, 就有 $|\sqrt[n]{a}-1|<\varepsilon$, 即
\[
  \lim_{n\to\infty}\sqrt[n]{a}=1.
\]

\noindent\textbf{证法 2}\quad 令 $\sqrt[n]{a}-1=h_n$, 则 $h_n>0$, 而且
\[
  a=(1+h_n)^n=1+nh_n+\cdots+(h_n)^n>nh_n,
\]
从而
\[
  0<h_n<\frac an,\qquad \forall n\in\mathbb{N}.
\]
这样, $\forall\varepsilon>0$, 要使 $|\sqrt[n]{a}-1|=h_n<\varepsilon$, 只要 $\dfrac an<\varepsilon$. 现在取 $N=\left[\dfrac a\varepsilon\right]$, 则只要 $n>N$, 就有 $|\sqrt[n]{a}-1|<\varepsilon$, 即 $\lim\limits_{n\to\infty}\sqrt[n]{a}=1$。

用定义来证明极限的过程, 实际上就是对于相对确定的正数 $\varepsilon$ 去找相应的正整数 $N$ 的过程. 换句话说就是, 你任给我一个正数 $\varepsilon$, 我去给你找一个相应的正整数 $N$, 在找 $N$ 的过程中, $\varepsilon$ 是相对固定的. 而找 $N$ 的过程实质就是解不等式的过程. 因此, 为了使所解的不等式尽可能的简单, 我们可以将所估计量做适当的放大 (参见例 2.1.7), 但不能放得太大, 要保证从中能够找到所需要的正整数 $N$。

下面我们来讨论一下发散序列. 如何用肯定的语气来表述一个序列是发散的呢? 从定义知, 一个序列 $\{x_n\}$ 是发散的等价于任何数 $a\in\mathbb{R}$ 都不是它的极限. 这样, 上述问题就转化为怎样用肯定的语气来表述 $\{x_n\}$ 不收敛于 $a$. 大家已经知道, $\{x_n\}$ 收敛于 $a$ 意味着 $a$ 的任意小的邻域的外面只有序列中的有限多项. 因此, 若 $\{x_n\}$ 不收敛于 $a$, 则存在 $a$ 的一个邻域 $U(a,\varepsilon_0)$, 使得在它之外必有 $\{x_n\}$ 的无穷多项. 于是, 若用肯定的语气, 发散序列可以叙述为:

设 $\{x_n\}$ 是一个序列. 若 $\forall a\in\mathbb{R}$, $\exists\varepsilon_0>0$, 对 $\forall N\in\mathbb{N}$, 总存在 $n_0>N$, 使得 $|x_{n_0}-a|>\varepsilon_0$, 则称 $\{x_n\}$ 为\textbf{发散序列}。

仔细观察上述的表述, 不难发现, 在发散序列的定义中, 只是将极限定义中的``$\forall$''换成了``$\exists$'', ``$\exists$''换成了``$\forall$''. 尽管这样, 我们还是希望读者能理解这样表述的实质, 而不是形式地加以记忆。

\noindent\textbf{例 2.1.9}\quad 证明如下定义的序列是发散的:
\[
  x_n=
  \begin{cases}
    \dfrac1n, & n=2k-1,\\
    n, & n=2k,
  \end{cases}
  \qquad k=1,2,\cdots .
\]

\noindent\textbf{证明}\quad 对任意给定的 $a\in\mathbb{R}$, 取 $\varepsilon_0=1$, 则对于任给的正整数 $N$, 选取 $n_0=2k_0>\max\{N,a+1\}$, 其中 $k_0$ 为正整数, 便有 $n_0>N$, 而且 $x_{n_0}-a=2k_0-a>1=\varepsilon_0$, 即 $a$ 不是 $\{x_n\}$ 的极限. 由 $a$ 的任意性知序列 $\{x_n\}$ 发散。

\subsection{无穷小量}

作为序列极限的特例, 我们引入无穷小量的概念。

\noindent\textbf{定义 2.1.2}\quad 设 $\{x_n\}$ 是一个序列. 若 $x_n\to0(n\to\infty)$, 则称序列 $\{x_n\}$ 为\textbf{无穷小量}, 记为 $x_n=o(1)(n\to\infty)$。

这里需特别指出的是, 无穷小量并不是一个很小的量, 而是极限为零的一个变量。

\noindent\textbf{例 2.1.10}\quad 证明序列 $\left\{\dfrac{a^n}{n!}\right\}$ 是无穷小量, 这里 $a>1$。

\noindent\textbf{证明}\quad 注意到
\[
  \frac{a^n}{n!}
  =
  \frac a1\cdot\frac a2\cdots\frac a{[a]}\cdot
  \frac a{[a]+1}\cdots\frac an
  < C\frac an
  \quad(n>[a]),
\]
其中 $C=\dfrac{a^{[a]}}{[a]!}$ 是常数, $\forall\varepsilon>0$, 要使
\[
  \left|\frac{a^n}{n!}-0\right|<\varepsilon,
\]
只需
\[
  n>\frac{Ca}{\varepsilon}.
\]
于是, 取 $N=\left[\dfrac{Ca}{\varepsilon}\right]+1$, 则只要 $n>N$, 就有
\[
  \left|\frac{a^n}{n!}-0\right|<\varepsilon,
\]
即 $\lim\limits_{n\to\infty}\dfrac{a^n}{n!}=0$, 从而序列 $\left\{\dfrac{a^n}{n!}\right\}$ 是无穷小量。

\noindent\textbf{例 2.1.11}\quad 证明序列 $\{x_n\}$ 是无穷小量, 其中
\[
  x_n=
  \begin{cases}
    \dfrac1{2^n}, & n=2k-1,\\[0.4em]
    \dfrac1{\sqrt n}, & n=2k,
  \end{cases}
  \qquad n=1,2,\cdots .
\]

\noindent\textbf{证明}\quad 因为对 $\forall n\in\mathbb{N}$, 有
\[
  \frac1{2^n}<\frac1{\sqrt n},
\]
所以, $\forall\varepsilon>0$, 取 $N=\left[\dfrac1{\varepsilon^2}\right]$, 则当 $n>N$ 时, 有
\[
  |x_n-0|\leqslant\frac1{\sqrt n}<\varepsilon,
\]
即 $\lim\limits_{n\to\infty}x_n=0$. 这就证明了所述序列 $\{x_n\}$ 是无穷小量。

利用极限的定义容易证明无穷小量有如下的基本性质:

\noindent\textbf{定理 2.1.1}\quad 设 $\{x_n\}$ 是一个序列。

(1) $\{x_n\}$ 是无穷小量的充分必要条件是 $\{|x_n|\}$ 是无穷小量;

(2) 若 $\{x_n\}$ 是无穷小量, $M$ 是一个常数, 则 $\{Mx_n\}$ 是无穷小量;

(3) $\lim\limits_{n\to\infty}x_n=a$ 的充分必要条件是 $\{x_n-a\}$ 是无穷小量。

证明留作练习。

\subsection{无穷大量}

对于序列 $x_n=(-1)^n$ 和 $x_n=(-1)^nn\ (n=1,2,\cdots)$ 来说, 虽然它们都不存在极限, 但这两者有着本质的区别: 前者中的 $x_n$ 随着 $n$ 的增加在 $-1$ 和 $1$ 这两点上跳来跳去; 而后者, 则随着 $n$ 的增加, 其绝对值目标一致地趋向 $+\infty$. 因此我们有如下定义:

\noindent\textbf{定义 2.1.3}\quad 设 $\{x_n\}$ 是一个序列. 若 $\forall M>0$, $\exists N$, 当 $n>N$ 时, 有 $x_n>M$, 则称 $\{x_n\}$ 为\textbf{正无穷大量} (有时也称 $\{x_n\}$ 的极限为 $+\infty$, 记为 $\lim\limits_{n\to\infty}x_n=+\infty$); 若 $\forall M>0$, $\exists N$, 当 $n>N$ 时, 有 $x_n<-M$, 则称 $\{x_n\}$ 为\textbf{负无穷大量} (有时也称 $\{x_n\}$ 的极限为 $-\infty$, 记为 $\lim\limits_{n\to\infty}x_n=-\infty$); 若 $\{|x_n|\}$ 是正无穷大量, 则称 $\{x_n\}$ 为\textbf{无穷大量} (有时也称 $\{x_n\}$ 的极限为 $\infty$, 记为 $\lim\limits_{n\to\infty}x_n=\infty$)。

显然, 正无穷大量和负无穷大量都是无穷大量. 具有有穷极限和无穷极限的序列有着本质的区别, 必须加以区别. 因此, 当序列有有穷极限 $a$ 时, 我们说它\textbf{收敛于} $a$; 当序列有无穷极限 $a$ 时, 我们说它\textbf{发散到} $+\infty$, $-\infty$ 或 $\infty$. 在本书的后续部分, 如果没有特别说明, 凡提到极限存在, 均指有穷极限的情形; 若包括无穷极限时, 则说其\textbf{广义极限存在}, 此时也说序列是\textbf{广义收敛的}。

\noindent\textbf{例 2.1.12}\quad 设
\[
  x_n=1+\frac12+\frac13+\cdots+\frac1n\quad(n=1,2,\cdots),
\]
证明 $\lim\limits_{n\to\infty}x_n=+\infty$。

\noindent\textbf{证明}\quad 注意到, 当 $n>2^k$ 时, 有
\[
\begin{aligned}
  x_n
  &=1+\frac12+\frac13+\frac14+\cdots+\frac1{2^{k-1}+1}
    +\cdots+\frac1{2^k}+\cdots+\frac1n\\
  &>1+\frac12+\frac14+\frac14+\frac18+\frac18+\frac18+\frac18
    +\cdots+\frac1{2^k}+\frac1{2^k}+\cdots+\frac1{2^k}\\
  &=1+\frac12+\frac12+\frac12+\cdots+\frac12\\
  &>\frac k2.
\end{aligned}
\]
故 $\forall M>0$, 欲使
\[
  x_n>M,
\]
只需
\[
  k>2M.
\]
这样, 只要取 $N=2^{[2M]+1}$, 则当 $n>N$ 时, 有
\[
  x_n>\frac{[2M]+1}{2}>\frac{2M}{2}=M,
\]
从而由极限定义知 $\lim\limits_{n\to\infty}x_n=+\infty$ 成立。

无穷大量和无穷小量之间有如下的关系:

\noindent\textbf{定理 2.1.2}\quad $\{x_n\}$ 是无穷小量的充分必要条件是 $\left\{\dfrac1{x_n}\right\}$ 是无穷大量, 这里假定对任意的正整数 $n$, 有 $x_n\ne0$。

\noindent\textbf{证明}\quad 先证必要性. 设 $\lim\limits_{n\to\infty}x_n=0$, 则 $\forall M>0$ (即 $\varepsilon=\dfrac1M>0$), $\exists N$, 当 $n>N$ 时, 有
\[
  |x_n|<\frac1M,
\]
从而有
\[
  \left|\frac1{x_n}\right|>M,
\]
即 $\lim\limits_{n\to\infty}\dfrac1{x_n}=\infty$。

再证充分性. 设 $\lim\limits_{n\to\infty}\dfrac1{x_n}=\infty$, 则 $\forall\varepsilon>0$ (即 $M=\dfrac1\varepsilon>0$), $\exists N$, 当 $n>N$ 时, 有
\[
  \left|\frac1{x_n}\right|>\frac1\varepsilon,
\]
从而有
\[
  |x_n|<\varepsilon,
\]
即 $\lim\limits_{n\to\infty}x_n=0$。

作为本节的结束, 我们再举一例。

\noindent\textbf{例 2.1.13}\quad 设 $\lim\limits_{n\to\infty}x_n=a$, 试证 $\lim\limits_{n\to\infty}\dfrac{x_1+\cdots+x_n}{n}=a$, 这里 $a$ 可以是有限实数, $+\infty$ 或 $-\infty$。

\noindent\textbf{证明}\quad 先证 $a$ 是有限实数的情形. $\forall\varepsilon>0$, 由于 $\lim\limits_{n\to\infty}x_n=a$, 故 $\exists N_1$, 使当 $n>N_1$ 时, 有
\[
  |x_n-a|<\frac\varepsilon2.
\]
对于取定的正整数 $N_1$, 由于
\[
  |x_1-a|+|x_2-a|+\cdots+|x_{N_1}-a|
\]
就是一个不随 $n$ 的变化而变化的常数, 因此有
\[
  \lim_{n\to\infty}
  \frac{|x_1-a|+|x_2-a|+\cdots+|x_{N_1}-a|}{n}=0.
\]
于是, 又存在正整数 $N_2>0$, 当 $n>N_2$ 时, 有
\[
  \frac{|x_1-a|+|x_2-a|+\cdots+|x_{N_1}-a|}{n}<\frac\varepsilon2.
\]
现在取 $N=\max\{N_1,N_2\}$, 则当 $n>N$ 时, 有
\[
\begin{aligned}
  \left|\frac{x_1+\cdots+x_n}{n}-a\right|
  &\leqslant
  \frac{|x_1-a|+|x_2-a|+\cdots+|x_n-a|}{n}\\
  &=
  \frac{|x_1-a|+\cdots+|x_{N_1}-a|}{n}
  +
  \frac{|x_{N_1+1}-a|+\cdots+|x_n-a|}{n}\\
  &\leqslant \frac\varepsilon2+\frac{n-N_1}{2n}\varepsilon<\varepsilon,
\end{aligned}
\]
从而有 $\lim\limits_{n\to\infty}\dfrac{x_1+\cdots+x_n}{n}=a$。

下证 $a=+\infty$ 的情形. $\forall M>0$, 由 $\lim\limits_{n\to\infty}x_n=+\infty$ 知, $\exists N_1$, 当 $n>N_1$ 时, 有 $x_n>2M$, 且 $x_1+x_2+\cdots+x_{N_1}>0$, 从而有
\[
  \frac{x_1+\cdots+x_n}{n}
  >
  \frac{x_{N_1+1}+\cdots+x_n}{n}
  >
  \frac{2(n-N_1)}nM.
\]
现在取 $N=2N_1$, 则当 $n>N$ 时, 有
\[
  \frac{n-N_1}{n}=1-\frac{N_1}{n}>\frac12,
\]
从而有
\[
  \frac{x_1+\cdots+x_n}{n}
  >
  \frac{2(n-N_1)}nM>M.
\]
这就证明了 $\lim\limits_{n\to\infty}\dfrac{x_1+\cdots+x_n}{n}=+\infty$。

对 $a=-\infty$ 的情形, 证明是类似的, 请读者自己给出。

\noindent\textbf{思考题}

(1) 若 $\lim\limits_{n\to\infty}x_n=\infty$, 是否有 $\lim\limits_{n\to\infty}\dfrac{x_1+\cdots+x_n}{n}=\infty$?

(2) 若 $\{x_n\}$ 发散, 是否有 $\left\{\dfrac{x_1+\cdots+x_n}{n}\right\}$ 亦发散?

\section{序列极限的性质}

在上一节中, 我们已经对极限有了初步的认识. 对于一个序列, 我们要讨论的主要问题是:
\[
\begin{array}{ll}
  \text{(1)} & \text{它是否存在极限?}\\
  \text{(2)} & \text{如果极限存在, 如何来求出这一极限?}
\end{array}
\]
对一些简单的序列, 我们很容易看出它是否收敛, 甚至能看出它的极限是什么. 但对一般的序列来说, 这是两个相当困难的问题. 可以这么说, 整个数学分析自始至终都在讨论各种极限的存在性以及求法等问题. 极限论可以看成是分析学与代数学的主要区别。

在这一节, 我们来进一步讨论序列极限的基本性质。

\noindent\textbf{定义 2.2.1}\quad 设 $\{x_n\}$ 是一个序列. 若 $\exists M>0$, 对 $\forall n$, 有 $|x_n|\leqslant M$ 成立, 则称 $\{x_n\}$ 是\textbf{有界的}。

显然, 以上定义等价于数集 $\{x_n\}$ 是一个有界集。

若一个序列 $\{x_n\}$ 是有界的, 则记为 $x_n=O(1)$ (读做大 $O$); 若存在 $M_2>M_1>0$ 和正整数 $N$, 使得当 $n>N$ 时, 有 $M_1<|x_n|<M_2$, 则以 $x_n=O_0(1)$ 表示之。

\noindent\textbf{定理 2.2.1}\quad (1) 改变一个序列 $\{x_n\}$ 的有限多项, 不改变其敛散性; 当 $\{x_n\}$ 收敛时, 则不改变其极限值。

(2) (\textbf{唯一性}) 收敛序列的极限是唯一的。

(3) (\textbf{有界性}) 收敛序列是有界的。

\noindent\textbf{证明}\quad (1) 由极限定义立即可得。

(2) 用反证法证之. 如果结论不真, 则存在收敛序列 $\{x_n\}$ 和实数 $b\ne a$, 使得 $\lim\limits_{n\to\infty}x_n=a$ 且 $\lim\limits_{n\to\infty}x_n=b$. 不失一般性, 我们不妨设 $a<b$. 现在取 $\varepsilon_0=\dfrac{b-a}{2}$, 则由极限定义知, 存在正整数 $N_1$ 和 $N_2$, 使得
\[
  |x_n-a|<\varepsilon_0,\qquad \forall n>N_1,
\]
\[
  |x_n-b|<\varepsilon_0,\qquad \forall n>N_2.
\]
令 $N=\max\{N_1,N_2\}$, 则当 $n>N$ 时, 上述两个不等式都成立. 由第一个不等式, 得
\[
  x_{N+1}<a+\varepsilon_0=\frac{a+b}{2},
\]
而由第二个不等式, 得
\[
  \frac{a+b}{2}=b-\varepsilon_0<x_{N+1}.
\]
这显然是不可能的, 因此 (2) 成立。

(3) 设 $\lim\limits_{n\to\infty}x_n=a$. 特别取 $\varepsilon_0=1$, 则 $\exists N$, 使得当 $n>N$ 时, 有 $|x_n-a|<1$, 从而
\[
  |x_n|\leqslant |a|+|x_n-a|<|a|+1,\qquad \forall n>N.
\]
令 $M=\max\{|a|+1,|x_1|,\cdots,|x_N|\}$, 则对 $\forall n$, 有 $|x_n|\leqslant M$, 即 $\{x_n\}$ 是有界的。

\noindent\textbf{思考题}\quad 试举例说明定理 2.2.1(3) 的逆命题不真。

\noindent\textbf{定理 2.2.2 (保序性)}\quad 给定两个序列 $\{x_n\}$ 和 $\{y_n\}$, 并且假定
\[
  \lim_{n\to\infty}x_n=a,\qquad
  \lim_{n\to\infty}y_n=b,
\]
则有:

(1) 若 $a<b$, 则对任意给定的 $c\in(a,b)$, $\exists N_0>0$, 使得当 $n>N_0$ 时, 有 $x_n<c<y_n$;

(2) 若 $\exists N_0>0$, 当 $n>N_0$ 时, 有 $x_n\leqslant y_n$, 则 $a\leqslant b$。

\noindent\textbf{证明}\quad (1) 令 $\varepsilon_0=\min\{c-a,b-c\}$, 则由极限的定义知, $\exists N>0$, 使得当 $n>N$ 时, 有 $|x_n-a|<\varepsilon_0$ 且 $|y_n-b|<\varepsilon_0$ 成立. 由此可得, 当 $n>N$ 时, 有
\[
  x_n<\varepsilon_0+a\leqslant(c-a)+a=c=b-(b-c)\leqslant b-\varepsilon_0<y_n.
\]

(2) 用反证法证之. 如若不然, 则有 $a>b$. 现取 $\varepsilon_0=\dfrac{a-b}{2}$, 由 (1) 所证知, $\exists N>0$, 使得当 $n>N$ 时, 有
\[
  y_n<\frac{a+b}{2}<x_n.
\]
这与已知条件 $x_n\leqslant y_n\ (n>N_0)$ 相矛盾, 故必有 $a\leqslant b$ 成立。

\noindent\textbf{推论}\quad 设 $\lim\limits_{n\to\infty}x_n=a>0$, 则对任给 $0<a'<a$, $\exists N_1>0$, 使得当 $n>N_1$ 时, 有 $x_n>a'$; 若 $\lim\limits_{n\to\infty}x_n=b<0$, 则对任给 $0>b'>b$, $\exists N_2>0$, 使得当 $n>N_2$ 时, 有 $x_n<b'$。

\noindent\textbf{注}\quad 注意, 在定理 2.2.2 的 (2) 中, 即使对一切的 $n$ 有 $x_n<y_n$, 也未必有 $a<b$. 例如,
\[
  x_n=1-\frac1n<1+\frac1n=y_n,\qquad n=1,2,\cdots,
\]
而
\[
  \lim_{n\to\infty}x_n=\lim_{n\to\infty}y_n=1.
\]

\noindent\textbf{定理 2.2.3 (极限的四则运算)}\quad 设 $\lim\limits_{n\to\infty}x_n=a$, $\lim\limits_{n\to\infty}y_n=b$, 则

(1) $\lim\limits_{n\to\infty}(x_n\pm y_n)=a\pm b$;

(2) $\lim\limits_{n\to\infty}(x_ny_n)=ab$;

(3) $\lim\limits_{n\to\infty}\dfrac{x_n}{y_n}=\dfrac ab$, 其中 $b\ne0,\ y_n\ne0$。

\noindent\textbf{证明}\quad (1) $\forall\varepsilon>0$, 由 $\lim\limits_{n\to\infty}x_n=a$ 知, $\exists N_1>0$, 使得当 $n>N_1$ 时, 有
\[
  |x_n-a|<\frac\varepsilon2;
\]
又由 $\lim\limits_{n\to\infty}y_n=b$ 知, $\exists N_2>0$, 使得当 $n>N_2$ 时, 有
\[
  |y_n-b|<\frac\varepsilon2.
\]
取 $N=\max\{N_1,N_2\}$, 则当 $n>N$ 时, 有
\[
  |(x_n\pm y_n)-(a\pm b)|
  \leqslant |x_n-a|+|y_n-b|
  <\frac\varepsilon2+\frac\varepsilon2=\varepsilon,
\]
即 $\lim\limits_{n\to\infty}(x_n\pm y_n)=a\pm b$。

(2) 首先, 我们有
\[
\begin{aligned}
  |x_ny_n-ab|
  &=|x_ny_n-y_na+y_na-ab|\\
  &\leqslant |y_n||x_n-a|+|a||y_n-b|.
\end{aligned}
\]
(这里, 为了与已知的极限发生联系, 我们使用了加一项、减一项的插项方法, 它是一个常用的方法)。

由有界性定理知, $\exists M_1>0$, $\forall n$, $|y_n|\leqslant M_1$. 令
\[
  M=\max\{M_1,|a|\}>0.
\]
$\forall\varepsilon>0$, 由 $\lim\limits_{n\to\infty}x_n=a$ 知, $\exists N_1>0$, 使得当 $n>N_1$ 时, 有
\[
  |x_n-a|<\frac{\varepsilon}{2M};
\]
又由 $\lim\limits_{n\to\infty}y_n=b$ 知, $\exists N_2>0$, 使得当 $n>N_2$ 时, 有
\[
  |y_n-b|<\frac{\varepsilon}{2M}.
\]
取 $N=\max\{N_1,N_2\}$, 则当 $n>N$ 时, 有
\[
\begin{aligned}
  |x_ny_n-ab|
  &\leqslant |y_n||x_n-a|+|a||y_n-b|\\
  &\leqslant M(|x_n-a|+|y_n-b|)\\
  &\leqslant M\left(\frac{\varepsilon}{2M}+\frac{\varepsilon}{2M}\right)\\
  &=\varepsilon,
\end{aligned}
\]
即 $\lim\limits_{n\to\infty}(x_ny_n)=ab$。

(3) 由 (2) 成立, 要证 (3), 只需证 $\lim\limits_{n\to\infty}\dfrac1{y_n}=\dfrac1b$ 即可. 首先, 我们有
\[
  \left|\frac1{y_n}-\frac1b\right|
  =
  \left|\frac{y_n-b}{y_nb}\right|.
\]
然后, 由 $\lim\limits_{n\to\infty}y_n=b\ne0$, 从定理 2.2.2 的推论知, $\exists N_1>0$, 使得当 $n>N_1$ 时, 有
\[
  |y_n|>\frac{|b|}{2}.
\]
这样, $\forall\varepsilon>0$, 由 $\lim\limits_{n\to\infty}y_n=b$ 知, $\exists N_2>0$, 使得当 $n>N_2$ 时, 有
\[
  |y_n-b|<\frac{|b|^2}{2}\varepsilon.
\]
取 $N=\max\{N_1,N_2\}$, 则当 $n>N$ 时, 有
\[
  \left|\frac1{y_n}-\frac1b\right|
  =
  \left|\frac{y_n-b}{y_nb}\right|
  <
  \frac2{|b|^2}\frac{|b|^2}{2}\varepsilon
  =
  \varepsilon,
\]
即 $\lim\limits_{n\to\infty}\dfrac1{y_n}=\dfrac1b$。

\noindent\textbf{例 2.2.1}\quad 求下列极限:

(1) $\lim\limits_{n\to\infty}\dfrac{3n^2+4n-100}{4n^2+5n+10^5}$;

(2) $\lim\limits_{n\to\infty}(1+q+q^2+\cdots+q^n)\ (|q|<1)$。

\noindent\textbf{解}\quad (1) 由于
\[
  \lim_{n\to\infty}\frac{3n^2+4n-100}{4n^2+5n+10^5}
  =
  \lim_{n\to\infty}
  \frac{3+\dfrac4n-\dfrac{100}{n^2}}
       {4+\dfrac5n+\dfrac{10^5}{n^2}},
\]
而
\[
  \lim_{n\to\infty}\left(3+\frac4n-\frac{100}{n^2}\right)=3,
\]
\[
  \lim_{n\to\infty}\left(4+\frac5n+\frac{10^5}{n^2}\right)=4,
\]
故有
\[
  \lim_{n\to\infty}\frac{3n^2+4n-100}{4n^2+5n+10^5}=\frac34.
\]

(2) 由于
\[
  1+q+q^2+\cdots+q^n=\frac{1-q^{n+1}}{1-q},
\]
而 $|q|<1$ 蕴涵着 $\lim\limits_{n\to\infty}q^n=0$, 于是
\[
  \lim_{n\to\infty}(1+q+q^2+\cdots+q^n)
  =
  \lim_{n\to\infty}\frac{1-q^{n+1}}{1-q}
  =
  \frac1{1-q}.
\]

\noindent\textbf{定理 2.2.4 (夹逼收敛原理)}\quad 设序列 $\{x_n\}$, $\{y_n\}$ 和 $\{z_n\}$ 满足
\[
  x_n\leqslant z_n\leqslant y_n,\qquad \forall n>N_0.
\]
若 $\lim\limits_{n\to\infty}x_n=\lim\limits_{n\to\infty}y_n=a$, 则 $\lim\limits_{n\to\infty}z_n=a$。

\noindent\textbf{证明}\quad $\forall\varepsilon>0$, 由 $\lim\limits_{n\to\infty}x_n=\lim\limits_{n\to\infty}y_n=a$ 知, 存在正整数 $N_1$ 和 $N_2$, 使得
\[
  |x_n-a|<\varepsilon,\qquad \forall n>N_1,
\]
\[
  |y_n-a|<\varepsilon,\qquad \forall n>N_2.
\]
令 $N=\max\{N_2,N_1,N_0\}$, 则当 $n>N$ 时, 有
\[
  a-\varepsilon<x_n\leqslant z_n\leqslant y_n<a+\varepsilon,
\]
即
\[
  |z_n-a|<\varepsilon.
\]
这表明 $\lim\limits_{n\to\infty}z_n=a$。

\noindent\textbf{例 2.2.2}\quad 求下列极限:

(1) $\lim\limits_{n\to\infty}\sqrt[n]{n}$;

(2) $\lim\limits_{n\to\infty}(a_1^n+a_2^n+\cdots+a_m^n)^{\frac1n}\ (a_i>0,\ i=1,2,\cdots,m)$。

\noindent\textbf{解}\quad (1) 令 $\sqrt[n]{n}=1+h_n$, 则 $h_n>0$, 而且
\[
  n=(1+h_n)^n
  =
  1+nh_n+\frac{n(n-1)}2h_n^2+\cdots+h_n^n
  >
  \frac{n(n-1)}2h_n^2\quad(n>3).
\]
于是
\[
  0<h_n<\frac{\sqrt2}{\sqrt{n-1}}.
\]
由于 $\lim\limits_{n\to\infty}\dfrac{\sqrt2}{\sqrt{n-1}}=0$, 由夹逼收敛原理知, $\lim\limits_{n\to\infty}h_n=0$, 从而 $\lim\limits_{n\to\infty}\sqrt[n]{n}=1$。

(2) 令 $a=\max\{a_1,a_2,\cdots,a_m\}$, 则有
\[
  a^n\leqslant a_1^n+a_2^n+\cdots+a_m^n\leqslant ma^n.
\]
于是
\[
  a\leqslant (a_1^n+a_2^n+\cdots+a_m^n)^{\frac1n}\leqslant \sqrt[n]{m}\,a.
\]
由于当 $n\geqslant m$ 时有
\[
  1\leqslant\sqrt[n]{m}\leqslant\sqrt[n]{n},
\]
由夹逼收敛原理和 (1) 知 $\lim\limits_{n\to\infty}\sqrt[n]{m}=1$. 再由夹逼收敛原理, 得
\[
  \lim_{n\to\infty}(a_1^n+a_2^n+\cdots+a_m^n)^{\frac1n}
  =
  a
  =
  \max\{a_1,a_2,\cdots,a_m\}.
\]

\noindent\textbf{例 2.2.3}\quad 设 $x_n\to0(n\to\infty)$, $p$ 为正整数. 证明 $\lim\limits_{n\to\infty}\sqrt[p]{1+x_n}=1$。

\noindent\textbf{证明}\quad 由题设, 不妨假定 $|x_n|<1$ 对一切 $n$ 成立. 这样, 我们有
\[
\begin{aligned}
  1-|x_n|
  &=\sqrt[p]{(1-|x_n|)^p}\\
  &\leqslant \sqrt[p]{1-|x_n|}
   \leqslant \sqrt[p]{1+x_n}
   \leqslant \sqrt[p]{1+|x_n|}\\
  &\leqslant \sqrt[p]{(1+|x_n|)^p}
  =1+|x_n|.
\end{aligned}
\]
而
\[
  \lim_{n\to\infty}(1-|x_n|)
  =
  \lim_{n\to\infty}(1+|x_n|)
  =
  1,
\]
所以, 由夹逼收敛原理知 $\lim\limits_{n\to\infty}\sqrt[p]{1+x_n}=1$ 成立。

\noindent\textbf{例 2.2.4}\quad 设
\[
  a_1=2,\qquad a_{n+1}=2+\frac1{a_n}\quad(n=1,2,\cdots),
\]
试问 $\{a_n\}$ 是否收敛? 若收敛则求其极限。

\noindent\textbf{分析}\quad 首先我们发现, 若 $\lim\limits_{n\to\infty}a_n=a$, 则必有
\[
  a=2+\frac1a.
\]
注意到 $a_n>0$, 从而得 $a=\sqrt2+1$。

令 $a_n=\sqrt2+1+h_n,\ (n=1,2,\cdots)$, 则 $\{a_n\}$ 收敛的充要条件是 $h_n=o(1)(n\to\infty)$。

\noindent\textbf{解}\quad 令 $a_n=\sqrt2+1+h_n,\ n=1,2,\cdots$. 我们有
\[
  1+\sqrt2+h_{n+1}=2+\frac1{1+\sqrt2+h_n},
\]
整理得
\[
  h_{n+1}=\frac{h_n(1-\sqrt2)}{1+\sqrt2+h_n}.
\]
显然 $0<|h_1|=\sqrt2-1<1$. 对 $k\geqslant1$, 若 $|h_k|<1$, 则有
\[
  |h_{k+1}|<|h_k|\frac{\sqrt2-1}{\sqrt2}<1.
\]
因此对所有的正整数 $n$, 有
\[
  |h_n|<1.
\]
再利用 $h_n$ 的递推关系我们推知
\[
  |h_{n+1}|\leqslant\frac12|h_n|,\qquad n=1,2,\cdots,
\]
从而对任意的 $n$, 有
\[
  0<|h_n|<\frac1{2^{n-1}}.
\]
由于 $\lim\limits_{n\to\infty}\dfrac1{2^{n-1}}=0$, 由夹逼收敛原理得 $\lim\limits_{n\to\infty}h_n=0$. 因此 $\{a_n\}$ 收敛, 并且 $\lim\limits_{n\to\infty}a_n=\sqrt2+1$。

作为本节的结束, 我们简要地讨论一下无穷小量和无穷大量的四则运算. 对任意两个无穷小量 $\{x_n\}$ 和 $\{y_n\}$, 由无穷小量定义和极限的四则运算定理知, 它们的和、差以及积仍然是无穷小量, 但它们的商可能有各种情形发生, 请看下例。

\noindent\textbf{例 2.2.5}\quad 设
\[
  x_n=\frac1{n+1},\qquad y_n=\frac1{n^2},\qquad z_n=\frac{(-1)^{n+1}}{n+1}\quad(n=1,2,\cdots).
\]
试考查:
\[
  \text{(1) }\lim_{n\to\infty}\frac{x_n}{y_n};\qquad
  \text{(2) }\lim_{n\to\infty}\frac{y_n}{x_n};\qquad
  \text{(3) }\lim_{n\to\infty}\frac{z_n}{x_n}.
\]

\noindent\textbf{解}\quad (1)
\[
  \lim_{n\to\infty}\frac{x_n}{y_n}
  =
  \lim_{n\to\infty}\frac{\dfrac1{n+1}}{\dfrac1{n^2}}
  =
  \lim_{n\to\infty}\frac{n^2}{n+1}
  =
  \infty;
\]
(2)
\[
  \lim_{n\to\infty}\frac{y_n}{x_n}
  =
  \lim_{n\to\infty}\frac{\dfrac1{n^2}}{\dfrac1{n+1}}
  =
  \lim_{n\to\infty}\frac{n+1}{n^2}
  =
  0;
\]
(3)
\[
  \lim_{n\to\infty}\frac{z_n}{x_n}
  =
  \lim_{n\to\infty}\frac{\dfrac{(-1)^{n+1}}{n+1}}{\dfrac1{n+1}}
  =
  \lim_{n\to\infty}(-1)^{n+1},
\]
极限不存在。

我们习惯把两个无穷小量的商形式地记为 $\dfrac00$, 上面例子告诉我们 $\dfrac00$ 是一个不定式, 即它的极限是不确定的, 可以有各种可能情况发生。

现在我们再来讨论无穷大量的情况. 我们有:
\[
\begin{array}{ll}
  \text{(1)} & (\pm\infty)+(\pm\infty)=\pm\infty;\\
  \text{(2)} & \infty\cdot\infty=\infty;\\
  \text{(3)} & \infty\cdot O_0(1)=\infty.
\end{array}
\]
但 $(+\infty)-(+\infty)$, $\dfrac{\infty}{\infty}$ 等则是不定式. 请读者作为练习, 举例来说明它们的不定性。

对于以上的形式记号, 它们的意义是明确的. 例如, (3) 中表示的意思是一个无穷大量与一个绝对值有正下界的有界序列的积仍然是一个无穷大量。

\noindent\textbf{例 2.2.6}\quad 设 $x_n=a_kn^k+a_{k-1}n^{k-1}+\cdots+a_0(a_k\ne0,k\geqslant1)$, 证明 $\lim\limits_{n\to\infty}x_n=\infty$。

\noindent\textbf{证明}\quad 设
\[
  x_n=n^k(a_k+a_{k-1}n^{-1}+a_{k-2}n^{-2}+\cdots+a_0n^{-k})=n^ky_n,
\]
由于
\[
  \lim_{n\to\infty}y_n=a_k\ne0,
\]
因此对所有充分大 $n$ 有
\[
  \frac{3|a_k|}{2}\geqslant |y_n|\geqslant\frac{|a_k|}{2},
\]
即 $y_n=O_0(1)$. 再由 $\lim\limits_{n\to\infty}n^k=\infty$, 即知
\[
  \lim_{n\to\infty}x_n=\infty.
\]

\noindent\textbf{思考题}\quad 对于这一例子试问下面的证明是否正确?
\[
  \lim_{n\to\infty}x_n
  =
  \lim_{n\to\infty}a_kn^k+\lim_{n\to\infty}a_{k-1}n^{k-1}+\cdots+\lim_{n\to\infty}a_0
  =
  \infty+\infty+\cdots+\infty+a_0=\infty.
\]

\section{单调收敛原理}

解决一个序列的收敛性问题, 只从定义出发往往是不够的, 还需要建立一系列理论结果. 本节我们就来介绍单调收敛原理和它的一些典型的应用范例。

\subsection{单调收敛原理}

若序列 $\{x_n\}$ 满足:
\[
  x_n\leqslant x_{n+1},\qquad \forall n\in\mathbb{N},
\]
则称 $\{x_n\}$ 是\textbf{单调递增} (上升) 的序列; 若满足:
\[
  x_n\geqslant x_{n+1},\qquad \forall n\in\mathbb{N},
\]
则称 $\{x_n\}$ 是\textbf{单调递减} (下降) 的序列; 单调上升的序列和单调下降的序列统称为\textbf{单调序列}. 细心的读者不难发现, 若将序列看成定义在 $\mathbb{N}$ 上的函数 $f(x)$, 则序列的单调性即是函数 $f(x)$ 的单调性。

\noindent\textbf{定理 2.3.1 (单调收敛原理)}\quad 单调有界序列必收敛。

\noindent\textbf{证明}\quad 这里只对单调上升的情形给出证明, 单调下降的情形请读者自己证明。

设 $\{x_n\}$ 单调上升, 而且有上界, 即
\[
  x_1\leqslant x_2\leqslant\cdots\leqslant x_n\leqslant x_{n+1}\leqslant\cdots\leqslant M,
\]
其中 $M$ 是一个正数. 现考虑集合 $E=\{x_n:n\in\mathbb{N}\}$, 则 $E$ 是一个非空有上界的集合, 故由确界定理, $E$ 必有上确界, 令 $a=\sup\{x_n\}$. 由上确界的定义知, $a$ 满足:

(1) $x_n\leqslant a,\ \forall n\in\mathbb{N}$;

(2) $\forall\varepsilon>0$, $\exists x_N$, 使得 $a-\varepsilon<x_N$。

这样, $\forall\varepsilon>0$, 当 $n>N$ 时, 有
\[
  a-\varepsilon<x_N\leqslant x_n\leqslant a,
\]
因此有
\[
  |x_n-a|<\varepsilon.
\]
这就证明了 $\lim\limits_{n\to\infty}x_n=a=\sup\{x_n\}$。

\noindent\textbf{注}\quad 容易证明: 若 $\{x_n\}$ 单调上升无上界, 则 $\lim\limits_{n\to\infty}x_n=+\infty$; 若 $\{x_n\}$ 单调下降无下界, 则 $\lim\limits_{n\to\infty}x_n=-\infty$. 这样, 结合定理 2.3.1, 我们知单调序列总是广义收敛的, 并且有

(1) 若 $\{x_n\}$ 单调上升, 则 $\lim\limits_{n\to\infty}x_n=\sup\{x_n\}$;

(2) 若 $\{x_n\}$ 单调下降, 则 $\lim\limits_{n\to\infty}x_n=\inf\{x_n\}$。

容易看出: 对一个单调递增序列 $\{x_n\}$, 若记 $A=\lim\limits_{n\to\infty}x_n$, 则或者存在 $N>0$ 使得 $x_n=x_N=A(\forall n>N)$, 或者 $x_n<A(\forall n>0)$。

\noindent\textbf{例 2.3.1}\quad 设 $A>0$, $x_1>0$, 定义
\[
  x_{n+1}=\frac12\left(x_n+\frac A{x_n}\right),\qquad n=1,2,\cdots.
\]
证明 $\lim\limits_{n\to\infty}x_n$ 存在, 并求出它的值。

\noindent\textbf{证明}\quad 易知这样产生的 $x_n$ 均为正数, 故有
\[
  x_{n+1}=\frac12\left(x_n+\frac A{x_n}\right)
  \geqslant \sqrt{x_n\frac A{x_n}}=\sqrt A
  \quad(n\geqslant2).
\]
这表明 $\{x_n\}$ 有下界. 此外,
\[
  x_{n+1}-x_n
  =
  \frac12\left(x_n+\frac A{x_n}\right)-x_n
  =
  \frac{A-x_n^2}{2x_n}\leqslant0
  \quad(n\geqslant2).
\]
这又证明了 $\{x_n\}$ 是单调下降的. 因此, $\lim\limits_{n\to\infty}x_n$ 存在, 设其为 $a$, 则在
\[
  x_{n+1}=\frac12\left(x_n+\frac A{x_n}\right)
\]
两边取极限, 得
\[
  a=\frac12\left(a+\frac Aa\right).
\]
解上述方程得 $a=\pm\sqrt A$. 注意到 $a\geqslant\sqrt A$, 便有 $a=\sqrt A$, 即 $\lim\limits_{n\to\infty}x_n=\sqrt A$。

\noindent\textbf{注}\quad 例 2.3.1 给出了一个求开方的近似方法. 例如, 令 $x_1=A=2$, 直接计算可得 $x_4=1.414257$, 这已经是 $\sqrt2$ 的一个很好的近似, 其实 $\sqrt2=1.414213\cdots$。

此外, 还需指出的是, 当我们知道了极限存在时, 可通过在序列的递推关系式中取极限而得到所求的极限值; 倘若我们并不知道极限是否存在, 则一般不能在递推式中取极限, 否则将得到荒谬的结果, 参见下面的例子。

\noindent\textbf{例 2.3.2}\quad 设 $x_1=1,\ x_{n+1}=-x_n,\ n=1,2,\cdots$, 易见 $\lim\limits_{n\to\infty}x_n$ 不存在. 但是, 若假定 $\lim\limits_{n\to\infty}x_n=a$, 并且在 $x_{n+1}=-x_n$ 两边取极限, 得 $a=-a$, 由此推出 $\lim\limits_{n\to\infty}x_n=0$ 的荒谬结果。

\noindent\textbf{例 2.3.3}\quad 设序列 $\{a_n\}$ 满足
\[
  a_1=1,\qquad a_{n+1}=1+\frac1{a_n},\qquad n=1,2,\cdots,
\]
证明 $\{a_n\}$ 收敛, 并求其极限。

\noindent\textbf{分析}\quad 直接计算, 得
\[
  a_1=1,\quad a_2=2,\quad a_3=\frac32,\quad a_4=\frac53,\quad a_5=\frac85.
\]
它们满足
\[
  a_1<a_3<a_5<a_4<a_2,
\]
由此我们推测: $\{a_{2k+1}\}$ 单调上升, 而 $\{a_{2k}\}$ 单调下降。

\noindent\textbf{证明}\quad 由 $a_n$ 满足的条件可导出
\begin{equation}
  a_{n+2}=2-\frac1{a_n+1}.
  \tag{2.3.1}
\end{equation}
于是
\[
\begin{aligned}
  a_{n+2}-a_n
  &=
  \left(2-\frac1{a_n+1}\right)
  -
  \left(2-\frac1{a_{n-2}+1}\right)\\
  &=
  \frac{a_n-a_{n-2}}{(a_{n-2}+1)(a_n+1)}.
\end{aligned}
\]
这表明 $a_{n+2}-a_n$ 与 $a_n-a_{n-2}$ 同号. 因此, 利用这一结果, 从 $a_1<a_3$ 和 $a_4<a_2$ 出发, 就可归纳地导出 $\{a_{2k+1}\}$ 单调上升和 $\{a_{2k}\}$ 单调下降. 此外, (2.3.1) 式和 $a_n$ 的递推公式蕴涵着
\[
  1<a_{2k},\qquad a_{2k+1}<2,\qquad \forall k\in\mathbb{N}.
\]
这样, 由单调收敛原理知, $\{a_{2k+1}\}$ 和 $\{a_{2k}\}$ 都收敛. 令
\[
  \lim_{k\to\infty}a_{2k+1}=a,\qquad
  \lim_{k\to\infty}a_{2k}=b,
\]
则有 $1<a,b<2$, 而且 (2.3.1) 式蕴涵着
\[
  a=2-\frac1{a+1},\qquad b=2-\frac1{b+1}.
\]
由此可导出
\[
  a=b=\frac12(1+\sqrt5).
\]
这就证明了 $\{a_n\}$ 收敛, 而且 $\lim\limits_{n\to\infty}a_n=\dfrac12(1+\sqrt5)$ (见本章习题 7)。

\subsection{无理数 $e$ 和欧拉常数 $c$}

作为单调收敛原理的应用, 我们给出无理数 $e$ 和欧拉常数 $c$ 的定义, 它们是数学中两个重要的常数。

考查序列
\[
  x_n=\left(1+\frac1n\right)^n,\qquad
  y_n=\left(1+\frac1n\right)^{n+1},\qquad n=1,2,\cdots.
\]
首先, 对任意的正整数 $n$, 有
\[
  y_n=\left(1+\frac1n\right)x_n>x_n.
\]
其次, 对任意的正整数 $n$, 有
\[
\begin{aligned}
  \frac{x_{n+1}}{x_n}
  &=
  \frac{(n+2)^{n+1}}{(n+1)^{n+1}}\frac{n^n}{(n+1)^n}\frac n{n+1}\\
  &=
  \left(1-\frac1{(n+1)^2}\right)^{n+1}\frac{n+1}{n}\\
  &>
  \left(1-\frac1{n+1}\right)\frac{n+1}{n}
  =
  1;
\end{aligned}
\]
\[
\begin{aligned}
  \frac{y_n}{y_{n+1}}
  &=
  \left(\frac{n+1}{n}\right)^{n+1}
  \left(\frac{n+1}{n+2}\right)^{n+2}\\
  &=
  \left(1+\frac1{n^2+2n}\right)^{n+1}\frac{n+1}{n+2}\\
  &>
  \left(1+\frac{n+1}{n^2+2n}\right)\frac{n+1}{n+2}\\
  &=
  \frac{n^3+4n^2+4n+1}{n^3+4n^2+4n}>1.
\end{aligned}
\]
这样我们就证明了: $\{x_n\}$ 单调上升有上界, $\{y_n\}$ 单调下降有下界. 因此, 它们二者都收敛, 而且有相同的极限. 我们约定用字母 $e$ 来表示其极限值, 即
\[
  \lim_{n\to\infty}\left(1+\frac1n\right)^n
  =
  \lim_{n\to\infty}\left(1+\frac1n\right)^{n+1}
  =
  e.
\]
$e$ 是数学中最重要的常数之一. 以后我们可以证明它是一个无理数并计算出具有任何给定精确度的近似值, 如
\[
  e=2.718281828459045\cdots.
\]
在数学的理论研究和应用中, 以 $e$ 为底的对数起着重要的作用, 这种对数称为\textbf{自然对数}, 记为 $\ln x$, 即
\[
  \ln x=\log_e x.
\]

此外, 上述讨论已经证明了如下的不等式
\begin{equation}
  \left(1+\frac1n\right)^n
  <
  e
  <
  \left(1+\frac1n\right)^{n+1},
  \qquad \forall n\in\mathbb{N}.
  \tag{2.3.2}
\end{equation}
对 (2.3.2) 两边取对数, 得
\[
  n\ln\left(1+\frac1n\right)
  <
  1
  <
  (n+1)\ln\left(1+\frac1n\right).
\]
由此即得
\begin{equation}
  \frac1{n+1}
  <
  \ln\left(1+\frac1n\right)
  <
  \frac1n,\qquad \forall n\in\mathbb{N}.
  \tag{2.3.3}
\end{equation}
从而
\[
  \sum_{k=1}^n\frac1{k+1}
  <
  \sum_{k=1}^n\ln\left(1+\frac1k\right)
  <
  \sum_{k=1}^n\frac1k.
\]
再注意到
\[
  \sum_{k=1}^n\ln\left(1+\frac1k\right)=\ln(n+1),
\]
便有
\begin{equation}
  \sum_{k=1}^n\frac1{k+1}
  <
  \ln(n+1)
  <
  \sum_{k=1}^n\frac1k,\qquad \forall n\in\mathbb{N}.
  \tag{2.3.4}
\end{equation}
现考虑序列
\[
  z_n=1+\frac12+\cdots+\frac1n-\ln n,\qquad n=1,2,\cdots.
\]
由不等式 (2.3.4), 可得
\[
  z_n>\ln(n+1)-\ln n>0
\]
和
\[
\begin{aligned}
  z_n-z_{n+1}
  &=
  \ln(n+1)-\ln n-\frac1{n+1}\\
  &=
  \ln\left(1+\frac1n\right)-\frac1{n+1}>0,
\end{aligned}
\]
其中最后一个不等式用到了不等式 (2.3.3). 这样, 我们就证明了 $\{z_n\}$ 单调下降有下界. 因此它收敛. 记
\[
  c=\lim_{n\to\infty}\left(1+\frac12+\cdots+\frac1n-\ln n\right),
\]
称 $c$ 为\textbf{欧拉 (Euler) 常数}. 我们可计算出
\[
  c=0.577216\cdots.
\]
欧拉常数也是一个重要的数学常数, 并且是一个不好研究的常数, 迄今为止人们还不知道它是有理数还是无理数。

由上面的讨论, 我们有
\[
  1+\frac12+\cdots+\frac1n=\ln n+c+\alpha_n,
\]
其中 $\alpha_n=o(1)(n\to\infty)$。

此外, 由不等式 (2.3.2), 可得
\[
  \frac{(n+1)^n}{n!}
  =
  \prod_{k=1}^n\left(1+\frac1k\right)^k
  <
  e^n
  <
  \prod_{k=1}^n\left(1+\frac1k\right)^{k+1}
  =
  \frac{(n+1)^{n+1}}{n!}.
\]
由此可得
\[
  \left(\frac{n+1}{e}\right)^n<n!<\left(\frac{n+1}{e}\right)^n(n+1),
\]
从而
\[
  \frac{n+1}{e}<\sqrt[n]{n!}<\frac{n+1}{e}\sqrt[n]{n+1}.
\]
这说明 $\lim\limits_{n\to\infty}\sqrt[n]{n!}=+\infty$, 而且
\[
  \lim_{n\to\infty}\frac n{\sqrt[n]{n!}}=e.
\]

这样, 我们对无穷大量 $\left\{1+\dfrac12+\cdots+\dfrac1n\right\}$ 和 $\{n!\}$ 的``量级''就有了一个比较清楚的认识, 前者是一个``弱得惊人''的无穷大量, 而后者是一个``大得出奇''的无穷大量。

\noindent\textbf{例 2.3.4}\quad 求极限
\[
  \lim_{n\to\infty}\left(1-\frac12+\frac13-\cdots+\frac{(-1)^{n-1}}n\right).
\]

\noindent\textbf{解}\quad 记
\[
  S_n=1-\frac12+\frac13-\cdots+\frac{(-1)^{n-1}}n,
\]
则有
\[
\begin{aligned}
  \lim_{n\to\infty}S_{2n}
  &=
  \lim_{n\to\infty}
  \left(1+\frac13+\cdots+\frac1{2n-1}
  -\left(\frac12+\cdots+\frac1{2n}\right)\right)\\
  &=
  \lim_{n\to\infty}
  \left(1+\frac12+\cdots+\frac1{2n}
  -2\left(\frac12+\cdots+\frac1{2n}\right)\right)\\
  &=
  \lim_{n\to\infty}
  \left(c+\ln2n+\alpha_{2n}-(c+\ln n+\alpha_n)\right)\\
  &=
  \ln2,
\end{aligned}
\]
其中 $c$ 是欧拉常数, $\alpha_n=o(1)(n\to\infty)$. 注意到
\[
  \lim_{n\to\infty}S_{2n+1}
  =
  \lim_{n\to\infty}\left(S_{2n}+\frac1{2n+1}\right)
  =
  \ln2,
\]
因此有 (见本章习题 7)
\[
  \lim_{n\to\infty}S_n=\ln2.
\]

\section{实数系连续性的基本定理}

这一节我们给出实数系连续性的几个常用的等价表述, 它们各有特点, 各自适合于不同的应用场合. 本节中的很多定理是许多数学论证中的基本工具。

\subsection{闭区间套定理}

\noindent\textbf{定理 2.4.1 (闭区间套定理)}\quad 设 $\{[a_n,b_n]\}$ 是一列闭区间, 并满足:

(1) $[a_n,b_n]\supseteq[a_{n+1},b_{n+1}],\ n=1,2,\cdots$;

(2) $\lim\limits_{n\to\infty}(b_n-a_n)=0$,

则存在唯一的一点 $c\in\mathbb{R}$, 使得 $c\in[a_n,b_n],\ n=1,2,\cdots$, 即
\[
  \{c\}=\bigcap_{n=1}^\infty [a_n,b_n].
\]

\noindent\textbf{证明}\quad 由条件 (1), 知
\[
  a_n\leqslant a_{n+1}\leqslant b_{n+1}\leqslant b_n,\qquad \forall n\in\mathbb{N}.
\]
这表明 $\{a_n\}$ 单调上升有上界, 并且 $\forall n\in\mathbb{N}$, $b_n$ 均是 $\{a_n\}$ 的一个上界; 另一方面, $\{b_n\}$ 单调下降有下界, 并且 $\forall n\in\mathbb{N}$, $a_n$ 均是 $\{b_n\}$ 的一个下界. 因此, $\{a_n\}$ 和 $\{b_n\}$ 都收敛. 令
\[
  \lim_{n\to\infty}a_n=a=\sup\{a_n\},\qquad
  \lim_{n\to\infty}b_n=b=\inf\{b_n\}.
\]
则有
\[
  a_n\leqslant a\leqslant b\leqslant b_n,\qquad \forall n\in\mathbb{N}.
\]
再由条件 (2), 知
\[
  a-b=\lim_{n\to\infty}a_n-\lim_{n\to\infty}b_n
  =
  \lim_{n\to\infty}(a_n-b_n)=0.
\]
记 $c=a=b$, 则有
\[
  a_n\leqslant\lim_{n\to\infty}a_n=c=\lim_{n\to\infty}b_n\leqslant b_n,\qquad \forall n\in\mathbb{N}.
\]
若另有 $d$ 满足
\[
  a_n\leqslant d\leqslant b_n,\qquad \forall n\in\mathbb{N},
\]
则由夹逼收敛原理, 有
\[
  d=\lim_{n\to\infty}a_n=\lim_{n\to\infty}b_n=c,
\]
故 $c$ 是所有区间 $[a_n,b_n]$ 的唯一公共点。

\noindent\textbf{注 1}\quad 闭区间的条件是重要的, 若区间是开的, 则定理的结论不一定成立. 例如, 对于开区间列 $\left(0,\dfrac1n\right),\ n=1,2,\cdots$, 显然满足定理的条件 (1) 和 (2), 但
\[
  \bigcap_{n=1}^\infty\left(0,\frac1n\right)=\varnothing.
\]
另一方面, 若开区间列 $\{(a_n,b_n)\}$ 满足条件

(i) $a_n<a_{n+1}<b_{n+1}<b_n,\ n=1,2,\cdots$;

(ii) $\lim\limits_{n\to\infty}(b_n-a_n)=0$,

则定理 2.4.1 的结论仍然成立。

\noindent\textbf{注 2}\quad 若定理 2.4.1 中条件 (2) 不满足, 即 $\lim\limits_{n\to\infty}(a_n-b_n)\ne0$, 则由定理的证明过程知, 此时 $a\ne b$, 而且容易导出, 此时有
\[
  \bigcap_{n=1}^\infty [a_n,b_n]=[a,b].
\]

\noindent\textbf{例 2.4.1}\quad 设 $\{[a_n,b_n]\}$ 是一列闭区间, 并且满足:
\[
  [a_i,b_i]\cap[a_j,b_j]\ne\varnothing,\qquad \forall i,j\in\mathbb{N},
\]
则存在一点 $c\in\mathbb{R}$, 使得 $c\in[a_n,b_n],\ n=1,2,\cdots$。

\noindent\textbf{证明}\quad 注意到
\[
  [a,b]\cap[c,d]\ne\varnothing\Rightarrow a\leqslant d,\ c\leqslant b,
\]
即知闭区间列所满足的条件蕴涵着任意的 $b_n$ 均是 $\{a_n\}$ 的上界. 于是, 令
\[
  c=\sup_{n\in\mathbb{N}}\{a_n\},
\]
则有
\[
  a_n\leqslant c\leqslant b_n,\qquad n=1,2,\cdots .
\]

\noindent\textbf{注}\quad 本例中的闭区间列并不一定是区间套, 只满足``两两有公共点''的条件, 同样推出它们有公共点的结论. 如若再加条件 $b_n-a_n\to0(n\to\infty)$, 则可得到与闭区间套定理完全相同的结论。

\subsection{有限覆盖定理}

设 $A$ 是 $\mathbb{R}$ 中的一个子集, $\{E_\lambda\}_{\lambda\in\Lambda}$ 是 $\mathbb{R}$ 中的一族子集组成的集合, 其中 $\Lambda$ 是一个指标集. 若 $A\subseteq\bigcup_{\lambda\in\Lambda}E_\lambda$, 则称 $\{E_\lambda\}_{\lambda\in\Lambda}$ 是 $A$ 的一个\textbf{覆盖}. 若 $\{E_\lambda\}_{\lambda\in\Lambda}$ 是 $A$ 的一个覆盖, 而且对每个 $\lambda\in\Lambda$, $E_\lambda$ 均是一个开区间, 则称 $\{E_\lambda\}_{\lambda\in\Lambda}$ 为 $A$ 的一个\textbf{开覆盖}; 若 $\{E_\lambda\}_{\lambda\in\Lambda}$ 是 $A$ 的一个覆盖, 而且 $\Lambda$ 的元素只有有限多个, 则称 $\{E_\lambda\}_{\lambda\in\Lambda}$ 是 $A$ 的一个\textbf{有限覆盖}。

\noindent\textbf{定理 2.4.2 (有限覆盖定理)}\quad 设 $[a,b]$ 是一个闭区间, $\{E_\lambda\}_{\lambda\in\Lambda}$ 是 $[a,b]$ 的任意一个开覆盖, 则必存在 $\{E_\lambda\}_{\lambda\in\Lambda}$ 一个子集构成 $[a,b]$ 的一个有限覆盖, 即在 $\{E_\lambda\}_{\lambda\in\Lambda}$ 中必有有限个开区间 $E_1,E_2,\cdots,E_N$, 使得 $[a,b]\subseteq\bigcup\limits_{j=1}^N E_j$。

\noindent\textbf{证明}\quad 如若不然, 我们将 $[a,b]$ 等分成两个闭区间, 则其中必有一个不能被 $\{E_\lambda\}_{\lambda\in\Lambda}$ 中的有限多个开区间所覆盖, 记其为 $[a_1,b_1]$; 再将 $[a_1,b_1]$ 等分成两个闭区间, 则必有其中的一个不能被 $\{E_\lambda\}_{\lambda\in\Lambda}$ 中的有限多个开区间所覆盖, 记其为 $[a_2,b_2]$. 如此进行下去, 就可得到一列闭区间 $[a_n,b_n]$, 满足:

(1) $[a_{n+1},b_{n+1}]\subseteq[a_n,b_n],\ \forall n$;

(2) $b_n-a_n=\dfrac{b-a}{2^n}\to0(n\to\infty)$;

(3) $\forall n$, $[a_n,b_n]$ 不能被 $\{E_\lambda\}_{\lambda\in\Lambda}$ 中的有限多个开区间所覆盖。

由闭区间套定理知, 存在唯一的 $\xi$, 使得 $\xi\in[a_n,b_n],\ \forall n\in\mathbb{N}$. 由于 $\xi\in[a,b]$, 而 $\{E_\lambda\}_{\lambda\in\Lambda}$ 是 $[a,b]$ 的一个开覆盖, 故必存在它的一个开区间 $E_{\lambda_0}=(c,d)$, 使得 $\xi\in(c,d)$, 即 $c<\xi<d$. 由于 $\{a_n\}$ 单调上升, $\{b_n\}$ 单调下降, 而且
\[
  \lim_{n\to\infty}a_n=\xi=\lim_{n\to\infty}b_n,
\]
故必有充分大的正整数 $n$, 使得
\[
  c<a_n\leqslant\xi\leqslant b_n<d,
\]
即 $[a_n,b_n]\subseteq E_{\lambda_0}$. 这与条件 (3) 矛盾. 定理得证。

关于有限覆盖定理, 我们必须注意以下两点:

(1) 不能将闭区间改为开区间. 如, $\left(\dfrac1n,\dfrac2n\right)(n=1,2,\cdots)$ 是 $(0,1)$ 的一个开覆盖, 但没有有限覆盖。

(2) 该定理是说, 对于闭区间的任意一个开覆盖必可从中找出有限多个开区间就可将该闭区间覆盖, 而并不是说, 任意一个闭区间都可被有限多个开区间覆盖. 若是这样, $(a-1,b+1)$ 总是 $[a,b]$ 的一个有限开覆盖。

在今后的应用中, 常可根据函数的局部性质得到一些开区间, 然后通过有限覆盖定理而得到函数在区间的整体性质。

\noindent\textbf{例 2.4.2}\quad 设函数 $f(x)$ 在 $[a,b]$ 上有定义, 且对任意的 $x\in[a,b]$, 都存在邻域 $U(x,\delta(x))(\delta(x)>0)$, 使得 $f(x)$ 在 $U(x,\delta(x))\cap[a,b]$ 上有界. 试证明 $f(x)$ 在 $[a,b]$ 上有界。

\noindent\textbf{证明}\quad 对任意的 $x\in[a,b]$, 由题设存在 $U(x,\delta(x))$, 使得 $f(x)$ 在 $U(x,\delta(x))\cap[a,b]$ 上有界, 记其界为 $M_x>0$, 即
\[
  |f(t)|\leqslant M_x,\qquad \forall t\in U(x,\delta(x))\cap[a,b].
\]
令
\[
  \mathcal{F}=\{U(x,\delta(x))=(x-\delta(x),x+\delta(x)):x\in[a,b]\}.
\]
显然 $\mathcal{F}$ 是 $[a,b]$ 的一个开覆盖. 由有限覆盖定理知, 存在 $\mathcal{F}$ 的有限个开区间
\[
  U(x_1,\delta(x_1)),\ U(x_2,\delta(x_2)),\ \cdots,\ U(x_N,\delta(x_N)),
\]
使得
\[
  [a,b]\subseteq\bigcup_{k=1}^N U(x_k,\delta(x_k)).
\]
令
\[
  M=\max\{M_{x_1},M_{x_2},\cdots,M_{x_N}\},
\]
则对任意的 $x\in[a,b]$, 必存在 $k(1\leqslant k\leqslant N)$, 使得 $x\in U(x_k,\delta(x_k))$, 从而
\[
  |f(x)|\leqslant M_{x_k}\leqslant M.
\]
这表明 $f(x)$ 在 $[a,b]$ 上有界。

\noindent\textbf{注}\quad 由上题的逆否命题我们可以得到以下有趣的事实: 若 $f(x)$ 在 $[a,b]$ 上无界, 则存在 $\xi\in[a,b]$, 使得 $\forall\delta>0$, $f(x)$ 在 $U(\xi,\delta)\cap[a,b]$ 上均无界。

\noindent\textbf{例 2.4.3}\quad 证明: 无论你用什么方法都不可能将 $[0,1]$ 区间中的全体实数排列成一个序列。

\noindent\textbf{证明}\quad 用反证法. 假定
\[
  [0,1]=\{x_1,x_2,x_3,\cdots,x_n,\cdots\}.
\]
取 $0<\varepsilon<\dfrac14$, 则开区间簇
\[
  \mathcal{F}=\left\{U\left(x_n,\frac{\varepsilon}{2^n}\right):n\in\mathbb{N}\right\}
\]
是 $[0,1]$ 区间的一个开覆盖. 由有限覆盖定理知, $\mathcal{F}$ 中存在有限个开区间
\[
  U\left(x_{n_1},\frac{\varepsilon}{2^{n_1}}\right),\
  U\left(x_{n_2},\frac{\varepsilon}{2^{n_2}}\right),\
  \cdots,\
  U\left(x_{n_\ell},\frac{\varepsilon}{2^{n_\ell}}\right),
\]
它们可将 $[0,1]$ 区间覆盖. 既然它们覆盖了 $[0,1]$ 区间, 那么这 $\ell$ 个区间的长度和必大于 $1$. 令
\[
  m=\max\{n_1,n_2,\cdots,n_\ell\},
\]
则有
\[
  1<
  2\sum_{k=1}^{\ell}\frac{\varepsilon}{2^{n_k}}
  <
  2\varepsilon\sum_{k=1}^m\frac1{2^k}
  =
  2\varepsilon\frac{1-\dfrac1{2^{m+1}}}{1-\dfrac12}
  <
  4\varepsilon<1.
\]
这一矛盾说明假设区间 $[0,1]$ 中的全体实数可以排列成一个序列是错误的。

\noindent\textbf{注}\quad 若一个数集 $E$ 中只有有限个元素或可以将它的所有元素排成一个序列, 则称 $E$ 是一个\textbf{可数集}. 上述例子说明区间 $[0,1]$ 不是可数集。

\subsection{聚点原理}

\noindent\textbf{定义 2.4.1}\quad 设 $E$ 是 $\mathbb{R}$ 中的一个子集. 若 $x_0\in\mathbb{R}$ ($x_0$ 不一定属于 $E$) 满足: 对 $\forall\delta>0$, 有 $U_0(x_0,\delta)\cap E\ne\varnothing$, 则称 $x_0$ 是 $E$ 的一个\textbf{聚点}。

\noindent\textbf{注}\quad (1) $x_0$ 是 $E$ 的聚点与 $x_0$ 是否属于 $E$ 无关;

(2) 由聚点的定义容易证明如下三个命题等价:

(i) $x_0$ 是 $E$ 的聚点;

(ii) $\forall\delta>0$, 在 $U(x_0,\delta)$ 中有 $E$ 的无穷多个点;

(iii) 存在 $E$ 中互异的点组成的序列 $\{x_n\}$, 使得 $\lim\limits_{n\to\infty}x_n=x_0$。

(3) 若 $x_0\in E$, 但它不是 $E$ 的聚点, 则称 $x_0$ 是 $E$ 的一个\textbf{孤立点}. 由定义, 此时必存在 $\delta>0$, 使得 $U(x_0,\delta)\cap E=\{x_0\}$。

\noindent\textbf{例 2.4.4}\quad 设
\[
  E=\left\{1,\frac12,\frac13,\cdots,\frac1n,\cdots\right\},
\]
则 $0$ 是它的唯一聚点, 而且 $0\notin E$; $\forall n\in\mathbb{N}$, $\dfrac1n$ 是 $E$ 的孤立点。

\noindent\textbf{例 2.4.5}\quad 设 $E$ 是 $[0,1]$ 中所有有理数组成的集合, 则 $E$ 的聚点全体是 $[0,1]$, 而 $E$ 没有孤立点. 注意, 此时 $E$ 是它的聚点集 $[0,1]$ 的真子集。

\noindent\textbf{定理 2.4.3 (聚点原理)}\quad $\mathbb{R}$ 中的任何一个有界无穷子集至少有一个聚点。

\noindent\textbf{证明}\quad 设 $E$ 是 $\mathbb{R}$ 中的一个有界无穷子集, 不妨设 $E\subseteq[a,b]$. 倘若定理的结论不真, 则 $\forall x\in[a,b]$, $x$ 都不是 $E$ 的聚点, 从而 $\exists\delta_x>0$, 使得 $U(x,\delta_x)\cap E$ 至多只有一个点. 显然,
\[
  \mathcal{F}=\{U(x,\delta_x):x\in[a,b]\}
\]
构成 $[a,b]$ 的一个开覆盖, 从而由有限覆盖定理知, 存在有限个开区间 $U(x_j,\delta_{x_j})(j=1,2,\cdots,m)$, 它们覆盖了 $[a,b]$, 即
\[
  [a,b]=\bigcup_{j=1}^m U(x_j,\delta_{x_j}).
\]
注意到 $E\subseteq[a,b]$, 而每个 $U(x_j,\delta_{x_j})$ 至多只有 $E$ 的一个点, 我们便可推知 $E$ 中至多有 $m$ 个点, 这与 $E$ 是 $\mathbb{R}$ 的无穷子集矛盾, 故 $E$ 必有聚点。

\noindent\textbf{思考题}\quad 试利用闭区间套定理证明聚点原理。

聚点原理与序列的子序列密切相关. 设 $\{x_n\}$ 是一个序列, 则由该序列的一部分元素按原来的顺序构成的序列 $\{x_{n_k}\}$ 称为是 $\{x_n\}$ 的一个\textbf{子序列}. 关于下标 $n_k$ 需要说明如下三点:

(1) $\{n_k\}$ 中的每一项都是正整数, 并且它是一个严格递增序列:
\[
  n_1<n_2<\cdots<n_k<\cdots;
\]

(2) $x_{n_k}$ 表示在子序列中它是第 $k$ 项, 在原序列中它是第 $n_k$ 项. 因此, 子序列的序号是 $k$, 而不是 $n_k$;

(3) 必有 $n_k\geqslant k$, 从而 $\lim\limits_{k\to\infty}n_k=+\infty$。

对于子序列, 我们容易证明如下的结果:

\noindent\textbf{定理 2.4.4}\quad 设 $\lim\limits_{n\to\infty}x_n=a$, 则对 $\{x_n\}$ 的任意子序列 $\{x_{n_k}\}$, 都有 $\lim\limits_{k\to\infty}x_{n_k}=a$。

此定理给出了判定一个序列发散的一个方法, 即: 若在一个序列 $\{x_n\}$ 中可以找到两个收敛的子序列 $\{x_{n_k}\}$ 和 $\{x_{n'_k}\}$, 使得
\[
  \lim_{k\to\infty}x_{n_k}=a\ne b=\lim_{k\to\infty}x_{n'_k},
\]
则该序列必然发散。

\noindent\textbf{例 2.4.6}\quad 设 $x_n=\sin\dfrac{n\pi}{2},\ n=1,2,\cdots$, 证明 $\{x_n\}$ 发散。

\noindent\textbf{证明}\quad 由于
\[
  \lim_{k\to\infty}x_{4k}
  =
  \lim_{k\to\infty}\sin2k\pi=0,
\]
而
\[
  \lim_{k\to\infty}x_{4k+1}
  =
  \lim_{k\to\infty}\sin\left(\frac\pi2+2k\pi\right)=1\ne0,
\]
因此 $\lim\limits_{n\to\infty}\sin\dfrac{n\pi}{2}$ 不存在。

\noindent\textbf{例 2.4.7}\quad 设
\[
  x_n=\left(1+\frac{(-1)^n}{n}\right)^n,\qquad n=1,2,\cdots,
\]
试证明 $\{x_n\}$ 为发散序列。

\noindent\textbf{证明}\quad 由于
\[
  \lim_{k\to+\infty}\left(1+\frac{(-1)^{2k}}{2k}\right)^{2k}=e,
\]
而
\[
  \lim_{k\to+\infty}\left(1+\frac{(-1)^{2k+1}}{2k+1}\right)^{2k+1}
  =
  \lim_{k\to+\infty}\frac1{\left(1+\dfrac1{2k}\right)^{2k}\left(1+\dfrac1{2k}\right)}
  =
  \frac1e,
\]
因此 $\lim\limits_{n\to\infty}x_n$ 不存在。

下面我们讲述波尔查诺-魏尔斯特拉斯 (Bolzano-Weierstrass) 定理。

\noindent\textbf{定理 2.4.5 (波尔查诺-魏尔斯特拉斯定理)}\quad 任何有界序列必有收敛的子序列。

\noindent\textbf{证明}\quad 设 $\{x_n\}$ 是一有界序列. 若 $\{x_n\}$ 只由有限多个数组成, 则它必有无穷多项等于同一个数, 此时定理结论自然成立。

现设 $\{x_n\}$ 由无穷多个互不相同的数组成, 则数集
\[
  E=\{x_n:n\in\mathbb{N}\}
\]
就是 $\mathbb{R}$ 中的一个有界无穷子集. 由聚点原理知, 它必有一个聚点, 设其为 $a$. 由聚点的定义知, 对任意的正整数 $k$, 在 $U(a,1/k)$ 中必有 $\{x_n\}$ 的无穷多项. 这样, 我们首先取
\[
  x_{n_1}\in U(a,1)\cap\{x_1,x_2,\cdots\};
\]
然后, 再取
\[
  x_{n_2}\in U\left(a,\frac12\right)\cap\{x_{n_1+1},x_{n_1+2},\cdots\}.
\]
如此进行下去, 就可找到 $\{x_n\}$ 的一个子列 $\{x_{n_k}\}$ 满足 $x_{n_k}\in U(a,1/k)$, 即
\[
  |x_{n_k}-a|<\frac1k,\qquad k=1,2,\cdots,
\]
从而 $\lim\limits_{k\to\infty}x_{n_k}=a$。

\subsection{柯西收敛准则}

一个序列是否收敛是我们研究的中心问题. 下面要介绍的柯西收敛准则从序列本身的性质出发给出了该序列收敛的充分必要条件. 在理论和应用上, 它都是十分重要的。

我们首先引入如下的定义。

\noindent\textbf{定义 2.4.2}\quad 设 $\{x_n\}$ 是一个序列, 若 $\forall\varepsilon>0$, $\exists N$, 当 $n,m>N$ 时, 有 $|x_n-x_m|<\varepsilon$, 则称 $\{x_n\}$ 是一个\textbf{柯西序列}。

\noindent\textbf{定理 2.4.6 (柯西收敛准则)}\quad 序列 $\{x_n\}$ 收敛的充分必要条件是它是一个柯西序列。

\noindent\textbf{证明}\quad \textbf{必要性}\quad 设序列 $\{x_n\}$ 收敛, 并设 $a=\lim\limits_{n\to\infty}x_n$, 则由极限的定义知, $\forall\varepsilon>0$, $\exists N$, 当 $n>N$ 时, 有
\[
  |x_n-a|<\frac\varepsilon2,
\]
从而当 $n,m>N$ 时, 有
\[
  |x_n-x_m|\leqslant |x_n-a|+|a-x_m|<\varepsilon.
\]

\textbf{充分性}\quad 设 $\{x_n\}$ 是一个柯西序列, 对 $\varepsilon_0=1$, 则 $\exists N$, 当 $n\geqslant N+1>N$ 时, 有
\[
  |x_n-x_{N+1}|<1,
\]
从而当 $n>N+1$ 时, 有 $|x_n|<1+|x_{N+1}|$. 令
\[
  M=\max\{|x_1|,|x_2|,\cdots,|x_N|,1+|x_{N+1}|\},
\]
则 $\forall n\in\mathbb{N}$, 有 $|x_n|\leqslant M$. 这表明 $\{x_n\}$ 是有界的. 于是由波尔查诺-魏尔斯特拉斯定理知, 存在 $\{x_n\}$ 的一个收敛子列 $\{x_{n_k}\}$, 记 $a=\lim\limits_{k\to\infty}x_{n_k}$。

下证
\[
  \lim_{n\to\infty}x_n=a.
\]
$\forall\varepsilon>0$, 由 $a=\lim\limits_{k\to\infty}x_{n_k}$ 知, $\exists K$, 当 $k\geqslant K$ 时, 有
\[
  |x_{n_k}-a|<\frac\varepsilon2.
\]
再由 $\{x_n\}$ 是柯西序列知, $\exists N$, 当 $n,m>N$ 时, 有
\[
  |x_n-x_m|<\frac\varepsilon2.
\]
这样, 取 $k>K$, 使得 $n_k>N$, 则当 $n>N$ 时, 有
\[
  |x_n-a|\leqslant |x_n-x_{n_k}|+|x_{n_k}-a|<\varepsilon,
\]
即 $\lim\limits_{n\to\infty}x_n=a$。

\noindent\textbf{注}\quad 我们有时候将一些数构成的集合 $K$ 称为是一个空间. 如果一个空间 $K$ 中的任何柯西序列 $\{x_n\}$ 都在 $K$ 中存在极限, 即存在 $A\in K$, 使得 $\lim\limits_{n\to\infty}x_n=A$, 则称 $K$ 是\textbf{完备的}. 例如, 上面证明的柯西收敛准则说明了实数域 $\mathbb{R}$ 是完备的. 显然, 有理数域 $\mathbb{Q}$ 是不完备的, 这是因为 $\mathbb{Q}$ 中的柯西序列的极限不一定是有理数. 另外, 任何开区间 $(a,b)$ 也是不完备的, 这是因为存在柯西序列 $\left\{b-\dfrac{b-a}{n}\right\}\subset(a,b)$ 但
\[
  \lim_{n\to\infty}\left(b-\frac{b-a}{n}\right)=b\notin(a,b).
\]

\noindent\textbf{例 2.4.8}\quad 证明序列 $\{x_n\}$ 是发散的, 其中
\[
  x_n=1+\frac12+\frac13+\cdots+\frac1n,\qquad n=1,2,\cdots.
\]

\noindent\textbf{证明}\quad 在本章的第一节我们已经用定义证明了该序列发散, 这里我们再用柯西收敛准则证明它发散. 取 $\varepsilon_0=\dfrac12$, 则对任意的正整数 $n$, 有
\[
  |x_{2n}-x_n|
  =
  \frac1{n+1}+\cdots+\frac1{2n}
  \geqslant \frac n{2n}=\frac12,
\]
所以它不是柯西序列, 从而它必是发散的。

\noindent\textbf{例 2.4.9}\quad 证明序列 $\{y_n\}$ 是收敛的, 其中
\[
  y_n=1+\frac1{2^2}+\frac1{3^2}+\cdots+\frac1{n^2},\qquad n=1,2,\cdots.
\]

\noindent\textbf{证明}\quad 由于对任意的 $n,m\in\mathbb{N}$, 我们不妨设 $n>m$, 有
\[
\begin{aligned}
  |y_n-y_m|
  &=
  \frac1{(m+1)^2}+\cdots+\frac1{n^2}\\
  &\leqslant
  \frac1{m(m+1)}+\cdots+\frac1{(n-1)n}
  \leqslant \frac1m.
\end{aligned}
\]
因此, $\forall\varepsilon>0$, 取 $N=\left[\dfrac1\varepsilon\right]$, 则当 $n,m>N$ 时, 就有
\[
  |y_n-y_m|<\varepsilon,
\]
即 $\{y_n\}$ 是柯西序列, 所以它是收敛的。

\noindent\textbf{例 2.4.10 (压缩映像原理)}\quad 设 $f(x)$ 在 $[a,b]$ 上有定义, $f([a,b])\subseteq[a,b]$, 且满足
\[
  |f(x)-f(y)|\leqslant q|x-y|,\qquad \forall x,y\in[a,b],
\]
其中 $0<q<1$. 证明: 存在唯一的 $c\in[a,b]$, 使得 $f(c)=c$。

\noindent\textbf{证明}\quad 任取 $x_0\in[a,b]$, 由条件 $f([a,b])\subseteq[a,b]$ 知, 我们可递推地定义
\[
  x_n=f(x_{n-1}),\qquad n=1,2,\cdots.
\]
由函数所满足的条件, 我们有
\[
  |x_{n+1}-x_n|
  =
  |f(x_n)-f(x_{n-1})|
  \leqslant q|x_n-x_{n-1}|,\qquad \forall n\in\mathbb{N}.
\]
反复应用上述不等式, 可得
\[
  |x_{n+1}-x_n|\leqslant q^n|x_1-x_0|,\qquad \forall n\in\mathbb{N}.
\]
因此, 对任意的正整数 $n$ 和 $p$, 有
\[
\begin{aligned}
  |x_{n+p}-x_n|
  &\leqslant \sum_{k=1}^p |x_{n+k}-x_{n+k-1}|\\
  &\leqslant \sum_{k=1}^p q^{n+k-1}|x_1-x_0|\\
  &=
  \frac{1-q^p}{1-q}q^n|x_1-x_0|\\
  &<
  \frac{q^n}{1-q}|x_1-x_0|=Lq^n,
\end{aligned}
\]
其中 $L=\dfrac1{1-q}|x_1-x_0|$. 由此立即可知 $\{x_n\}$ 是柯西序列, 从而它收敛, 记 $\lim\limits_{n\to\infty}x_n=c$. 显然, $c\in[a,b]$. 又由于
\[
  |f(x_n)-f(c)|\leqslant q|x_{n-1}-c|,
\]
故有 $\lim\limits_{n\to\infty}f(x_n)=f(c)$. 这样, 在等式 $x_{n+1}=f(x_n)$ 两边取极限, 即得 $c=f(c)$, 从而 $c$ 的存在性得证。

下证唯一性. 若存在 $c'\in[a,b]$ 也满足 $f(c')=c'$, 而且 $c'\ne c$, 则由
\[
  |c-c'|=|f(c)-f(c')|\leqslant q|c-c'|
\]
即得矛盾, 从而 $c$ 的唯一性得证。

\section{序列的上、下极限}

上极限和下极限的概念是极限概念的拓广, 它为研究序列的敛散性开拓了一条全新的思路。

设 $\{x_n\}$ 是有界序列. 令
\[
  \ell_n=\inf\{x_n,x_{n+1},x_{n+2},\cdots\},
  \qquad
  h_n=\sup\{x_n,x_{n+1},x_{n+2},\cdots\},
\]
则有
\[
  \ell_1\leqslant\ell_2\leqslant\cdots\leqslant\ell_n\leqslant\cdots\leqslant h_n\leqslant\cdots\leqslant h_2\leqslant h_1,
\]

即 $\{\ell_n\}$ 和 $\{h_n\}$ 是单调有界序列. 令
\[
  \ell=\sup_n\{\ell_n\}=\sup_n\inf_k\{x_{n+k}\},
  \qquad
  h=\inf_n\{h_n\}=\inf_n\sup_k\{x_{n+k}\},
\]
则 $\ell$ 与 $h$ 分别称为 $\{x_n\}$ 的\textbf{下极限}和\textbf{上极限}, 记为
\[
  \ell=\underline{\lim_{n\to\infty}}\,x_n,\qquad
  h=\overline{\lim_{n\to\infty}}\,x_n;
\]
有时下极限和上极限也记为 $\liminf\limits_{n\to\infty}x_n$ 和 $\limsup\limits_{n\to\infty}x_n$。

此外, 对于无界序列, 我们规定:

(1) 若 $\{x_n\}$ 无上界, 则 $\overline{\lim\limits_{n\to\infty}}\,x_n=+\infty$. 注意此时必有: $h_n=+\infty,\ \forall n\in\mathbb{N}$。

(2) 若 $\{x_n\}$ 无下界, 则 $\underline{\lim\limits_{n\to\infty}}\,x_n=-\infty$. 注意此时必有: $\ell_n=-\infty,\ \forall n\in\mathbb{N}$。

(3) 若 $\{x_n\}$ 有下界但无上界, 则 $\{\ell_n\}$ 有定义, 故可定义
\[
  \underline{\lim_{n\to\infty}}\,x_n=\sup_n\{\ell_n\},
\]
但须注意的是, 此时 $\sup_n\{\ell_n\}$ 可能为 $+\infty$。

(4) 若 $\{x_n\}$ 有上界但无下界, 则 $\{h_n\}$ 有定义, 故可定义
\[
  \overline{\lim_{n\to\infty}}\,x_n=\inf_n\{h_n\},
\]
但须注意的是, 此时 $\inf_n\{h_n\}$ 可能为 $-\infty$。

这样, 对任意的序列 $\{x_n\}$, $\overline{\lim\limits_{n\to\infty}}\,x_n$ 和 $\underline{\lim\limits_{n\to\infty}}\,x_n$ 都有明确的定义, 而且恒有
\[
  \underline{\lim_{n\to\infty}}\,x_n
  \leqslant
  \overline{\lim_{n\to\infty}}\,x_n.
\]

\noindent\textbf{例 2.5.1}\quad 试求序列 $\{x_n\}$ 的上极限与下极限, 其中
\[
  (1)\ x_n=\sin\frac n2\pi;\qquad
  (2)\ x_n=n^{(-1)^n};
\]
\[
  (3)\ x_n=(-1)^n n;\qquad
  (4)\ x_n=n.
\]

\noindent\textbf{解}\quad (1) 由 $\forall n\in\mathbb{N}$, $\ell_n=-1,\ h_n=1$, 得
\[
  \underline{\lim_{n\to\infty}}\,\sin\frac n2\pi=-1,
  \qquad
  \overline{\lim_{n\to\infty}}\,\sin\frac n2\pi=1.
\]

(2) 由 $\forall n\in\mathbb{N}$, $\ell_n=0,\ h_n=+\infty$, 即 $\{x_n\}$ 无上界, 得
\[
  \underline{\lim_{n\to\infty}}\, n^{(-1)^n}=0,
  \qquad
  \overline{\lim_{n\to\infty}}\, n^{(-1)^n}=+\infty.
\]

(3) 由 $\forall n\in\mathbb{N}$, $\ell_n=-\infty,\ h_n=+\infty$, 即 $\{x_n\}$ 既无上界又无下界, 得
\[
  \underline{\lim_{n\to\infty}}\,(-1)^n n=-\infty,
  \qquad
  \overline{\lim_{n\to\infty}}\,(-1)^n n=+\infty.
\]

(4) 由 $\forall n\in\mathbb{N}$, $\ell_n=n,\ h_n=+\infty$, 得
\[
  \underline{\lim_{n\to\infty}}\, n=+\infty,
  \qquad
  \overline{\lim_{n\to\infty}}\, n=+\infty.
\]

从上极限的定义, 容易证明如下关于上极限的几个等价的描述。

\noindent\textbf{定理 2.5.1}\quad 设 $\{x_n\}$ 是一有界序列, $h$ 是一实数, 则下列三个命题等价:

(1) $h$ 是 $\{x_n\}$ 的上极限;

(2) $\forall\varepsilon>0,\ \exists N$, 当 $n>N$ 时, 有 $x_n<h+\varepsilon$, 而且对 $\forall K,\ \exists n_k>K$, 使得 $x_{n_k}>h-\varepsilon$;

(3) 存在子列 $\{x_{n_k}\}$, 使得 $\lim\limits_{k\to\infty}x_{n_k}=h$, 而且对任何其他收敛子列 $\{x_{n'_k}\}$, 有 $\lim\limits_{k\to\infty}x_{n'_k}\leqslant h$。

\noindent\textbf{证明}\quad (1) $\Rightarrow$ (2). $\forall n\in\mathbb{N}$, 记 $h_n=\sup\{x_n,x_{n+1},\cdots\}$. 由 $\overline{\lim\limits_{n\to\infty}}\,x_n=h,\ \forall\varepsilon>0,\ \exists N$, 当 $n>N$ 时, 有
\[
  |h_n-h|<\varepsilon,
  \qquad\text{即}\qquad h-\varepsilon<h_n<h+\varepsilon.
\]
由于对任意的 $n\in\mathbb{N}$, 有 $x_n\leqslant h_n$, 因此, 当 $n>N$ 时, 有
\[
  x_n<h+\varepsilon.
\]
对 $\forall K>0$, 当 $n'_k>\max\{N,K\}$ 时, 有
\[
  h_{n'_k}>h-\varepsilon.
\]
由 $h_{n'_k}=\sup\{x_{n'_k},x_{n'_k+1},x_{n'_k+2},\cdots\}$ 知, 存在 $n_k\geqslant n'_k>K$ 使得有
\[
  x_{n_k}>h-\varepsilon.
\]
综上所述, (2) 成立。

(2) $\Rightarrow$ (3). 由 (2), 对 $\varepsilon=1$, 显然存在正整数 $n_1\geqslant1$, 使得
\[
  h-1<x_{n_1}<h+1.
\]
现假定对 $\varepsilon_k=\dfrac1k(k\geqslant1,\ k\in\mathbb{N})$, 存在 $n_k\in\mathbb{N}$ 满足
\[
  h-\frac1k<x_{n_k}<h+\frac1k.
\]
在 (2) 中取 $\varepsilon=\dfrac1{k+1}$, 从而存在 $N_{k+1}$, 当 $n>N_{k+1}$ 时, 有
\[
  x_n<h+\frac1{k+1}.
\]
令 $K=\max\{N_{k+1},n_k\}$. 由 (2) 知存在 $n_{k+1}>K$, 使得
\[
  h-\frac1{k+1}<x_{n_{k+1}}<h+\frac1{k+1}.
\]
由上述归纳法找到的 $\{x_{n_k}\}$ 显然满足 $\lim\limits_{k\to\infty}x_{n_k}=h$。

现设 $\{x_{n'_k}\}$ 是 $\{x_n\}$ 的一个收敛子列, 并设
\[
  \lim_{k\to\infty}x_{n'_k}=h'.
\]
由 (2) 知, $\forall\varepsilon>0,\ \exists N$, 当 $n>N$ 时, 有
\[
  x_n<h+\varepsilon.
\]
注意到 $n'_k\geqslant k$, 从而当 $k>N$ 时, 有
\[
  x_{n'_k}<h+\varepsilon.
\]
这说明
\[
  h'=\lim_{k\to\infty}x_{n'_k}\leqslant h+\varepsilon.
\]
由 $\varepsilon$ 的任意性, 知 $h'\leqslant h$. 因此 (3) 的结论成立。

(3) $\Rightarrow$ (1). 设 $\overline{\lim\limits_{n\to\infty}}\,x_n=h'$, 我们将证明 $h'=h$, 其中 $h$ 满足 (3) 中的条件。

由 (3) 知, 存在 $\{x_n\}$ 的子列 $\{x_{n_k}\}$ 满足:
\[
  \lim_{k\to\infty}x_{n_k}=h.
\]
注意到 $\forall k\in\mathbb{N}$, 有 $h_{n_k}\geqslant x_{n_k}$, 从而有 $h'\geqslant h$。

假若 $h'>h$, 我们将推出与 (3) 相矛盾的结果. 事实上, 令 $\varepsilon_0=\dfrac{h'-h}{2}>0$. 由 (2) 知存在 $\{x_n\}$ 的子列 $\{x_{n'_k}\}$, 使得对 $\forall k\in\mathbb{N}$, 有
\[
  x_{n'_k}>h'-\varepsilon_0=h+\varepsilon_0.
\]
由于 $\{x_n\}$ 为有界序列, 从而 $\{x_{n'_k}\}$ 存在收敛子列 $\{x_{n''_k}\}$, 显然有 $\lim\limits_{k\to\infty}x_{n''_k}\geqslant h+\varepsilon_0>h$. 这与 (3) 矛盾, 从而 $h'=h$, 即 (1) 成立。

\noindent\textbf{注}\quad 对于下极限, 我们有相应于定理 2.5.1 的结果, 请读者自己给出这些结果的叙述与说明。

\noindent\textbf{定理 2.5.2}\quad (1) 若有界序列 $\{x_n\}$ 由互不相同的数组成, 则上极限 $\overline{\lim\limits_{n\to\infty}}\,x_n$ 是 $\{x_n\}$ 的最大聚点, 而下极限 $\underline{\lim\limits_{n\to\infty}}\,x_n$ 是 $\{x_n\}$ 的最小聚点;

(2) $\{x_{n_k}\}$ 是 $\{x_n\}$ 的任一子列, 则
\[
  \underline{\lim_{n\to\infty}}\,x_n
  \leqslant
  \underline{\lim_{k\to\infty}}\,x_{n_k}
  \leqslant
  \overline{\lim_{k\to\infty}}\,x_{n_k}
  \leqslant
  \overline{\lim_{n\to\infty}}\,x_n;
\]

(3) $\lim\limits_{n\to\infty}x_n=a$ 的充分必要条件是 $\underline{\lim\limits_{n\to\infty}}\,x_n=\overline{\lim\limits_{n\to\infty}}\,x_n=a$, 其中 $a$ 可以是有限数、$-\infty$ 或 $+\infty$。

\noindent\textbf{证明}\quad (1) 设 $\overline{\lim\limits_{n\to\infty}}\,x_n=h$. 由定理 2.5.1(3) 知, 存在 $\{x_n\}$ 的子列 $\{x_{n_k}\}$, 使得 $\lim\limits_{k\to\infty}x_{n_k}=h$. 由于 $\{x_{n_k}\}$ 是由互不相同的数组成, $h$ 是 $\{x_{n_k}\}$ 的聚点, 因此 $h$ 更是 $\{x_n\}$ 的聚点. 对于 $\{x_n\}$ 的任何聚点 $h'$, 都存在它的子列 $\{x_{n'_k}\}$, 使得 $\lim\limits_{k\to\infty}x_{n'_k}=h'$. 由定理 2.5.1(3) 知必有 $h'\leqslant h$. 这就证明了 $h$ 是 $\{x_n\}$ 的最大聚点. 同样的方法可证 $\underline{\lim\limits_{n\to\infty}}\,x_n$ 是 $\{x_n\}$ 的最小聚点。

(2) 对 $\forall k\in\mathbb{N}$, 有
\[
  \ell_{n_k}=\inf\{x_{n_k},x_{n_k+1},\cdots\}
  \leqslant
  \inf\{x_{n_k},x_{n_{k+1}},\cdots\}.
\]
因此
\[
  \underline{\lim_{n\to\infty}}\,x_n
  =
  \lim_{n\to\infty}\ell_n
  =
  \lim_{k\to\infty}\ell_{n_k}
  \leqslant
  \lim_{k\to\infty}\inf\{x_{n_k},x_{n_{k+1}},\cdots\}
  =
  \underline{\lim_{k\to\infty}}\,x_{n_k}.
\]
对 $\forall k\in\mathbb{N}$, 有
\[
  h'_k=\sup\{x_{n_k},x_{n_{k+1}},\cdots\}
  \leqslant
  \sup\{x_{n_k},x_{n_k+1},\cdots\}=h_{n_k},
\]
因此
\[
  \overline{\lim_{k\to\infty}}\,x_{n_k}
  =
  \lim_{k\to\infty}h'_k
  \leqslant
  \lim_{k\to\infty}h_{n_k}
  =
  \lim_{n\to\infty}h_n
  =
  \overline{\lim_{n\to\infty}}\,x_n.
\]
(3) 我们只证 $a$ 为有限数的情形, 其余的证明留给读者作为练习。

充分性\quad 假定 $\underline{\lim\limits_{n\to\infty}}\,x_n=\overline{\lim\limits_{n\to\infty}}\,x_n=a$, 由于 $\forall n\in\mathbb{N}$, 我们有
\[
  \ell_n=\inf_k\{x_{n+k}\}\leqslant x_n\leqslant\sup_k\{x_{n+k}\}=h_n,
\]
由夹逼收敛原理知, $\lim\limits_{n\to\infty}x_n=a$ 成立。

必要性\quad 设 $\lim\limits_{n\to\infty}x_n=a$, 则 $\forall\varepsilon>0,\ \exists N$, 当 $n>N$ 时, 有
\[
  a-\varepsilon<x_n<a+\varepsilon,
\]
从而有
\[
  a-\varepsilon\leqslant\ell_n\leqslant x_n\leqslant h_n\leqslant a+\varepsilon,\qquad \forall n>N.
\]
这说明
\[
  \underline{\lim_{n\to\infty}}\,x_n
  =
  \lim_{n\to\infty}\ell_n
  =
  a
  =
  \lim_{n\to\infty}h_n
  =
  \overline{\lim_{n\to\infty}}\,x_n.
\]

\noindent\textbf{例 2.5.2}\quad 设 $\{x_n\}$ 是 $[0,1]$ 中全体有理数构成的序列, 试求它的上极限与下极限。

\noindent\textbf{解法 1}\quad 由于 $\{x_n\}$ 在 $[0,1]$ 稠密, 因此, 对任意的 $n\in\mathbb{N}$, $\{x_n,x_{n+1},x_{n+2},\cdots\}$ 仍在 $[0,1]$ 稠密. 由此, $\forall n\in\mathbb{N}$, 有 $\ell_n=0,\ h_n=1$, 从而有
\[
  \underline{\lim_{n\to\infty}}\,x_n=0,
  \qquad
  \overline{\lim_{n\to\infty}}\,x_n=1.
\]

\noindent\textbf{解法 2}\quad 由于 $\lim\limits_{k\to\infty}\dfrac1k=0,\ \lim\limits_{k\to\infty}\dfrac{k-1}{k}=1$, 因此 $0$ 与 $1$ 都是 $\{x_n\}$ 的聚点. 显然 $0$ 是它的最小聚点, $1$ 是它的最大聚点. 这说明: $\underline{\lim\limits_{n\to\infty}}\,x_n=0,\ \overline{\lim\limits_{n\to\infty}}\,x_n=1$。

下面的定理给出了上、下极限的保序性和运算性质。

\noindent\textbf{定理 2.5.3}\quad 设 $\{x_n\}$ 和 $\{y_n\}$ 是任意给定的两个有界序列, 则有

(1) $x_n\leqslant y_n(n=1,2,\cdots)$ 蕴涵着
\[
  \underline{\lim_{n\to\infty}}\,x_n
  \leqslant
  \underline{\lim_{n\to\infty}}\,y_n,
  \qquad
  \overline{\lim_{n\to\infty}}\,x_n
  \leqslant
  \overline{\lim_{n\to\infty}}\,y_n;
\]

(2)
\[
  \underline{\lim_{n\to\infty}}\,(-x_n)
  =
  -\overline{\lim_{n\to\infty}}\,x_n,
  \qquad
  \overline{\lim_{n\to\infty}}\,(-x_n)
  =
  -\underline{\lim_{n\to\infty}}\,x_n;
\]

(3)
\[
  \underline{\lim_{n\to\infty}}\,x_n
  +
  \underline{\lim_{n\to\infty}}\,y_n
  \leqslant
  \underline{\lim_{n\to\infty}}\,(x_n+y_n)
  \leqslant
  \underline{\lim_{n\to\infty}}\,x_n
  +
  \overline{\lim_{n\to\infty}}\,y_n,
\]
\[
  \overline{\lim_{n\to\infty}}\,x_n
  +
  \underline{\lim_{n\to\infty}}\,y_n
  \leqslant
  \overline{\lim_{n\to\infty}}\,(x_n+y_n)
  \leqslant
  \overline{\lim_{n\to\infty}}\,x_n
  +
  \overline{\lim_{n\to\infty}}\,y_n;
\]

(4) 若 $x_n\geqslant0,\ y_n\geqslant0,\ n=1,2,\cdots$, 则
\[
  \underline{\lim_{n\to\infty}}\,x_n\cdot
  \underline{\lim_{n\to\infty}}\,y_n
  \leqslant
  \underline{\lim_{n\to\infty}}\,(x_n\cdot y_n)
  \leqslant
  \underline{\lim_{n\to\infty}}\,x_n\cdot
  \overline{\lim_{n\to\infty}}\,y_n,
\]
\[
  \overline{\lim_{n\to\infty}}\,x_n\cdot
  \underline{\lim_{n\to\infty}}\,y_n
  \leqslant
  \overline{\lim_{n\to\infty}}\,(x_n\cdot y_n)
  \leqslant
  \overline{\lim_{n\to\infty}}\,x_n\cdot
  \overline{\lim_{n\to\infty}}\,y_n.
\]

\noindent\textbf{证明}\quad 我们只证明 (3) 中的前两个不等式, 其余的不等式留给读者作为练习。

取 $\{x_n+y_n\}$ 的一个收敛子列 $\{x_{n_k}+y_{n_k}\}$, 使得
\[
  \lim_{k\to\infty}(x_{n_k}+y_{n_k})
  =
  \underline{\lim_{n\to\infty}}\,(x_n+y_n).
\]
记 $\ell'_n=\inf\{x_n,x_{n+1},\cdots\},\ \ell''_n=\inf\{y_n,y_{n+1},\cdots\}$, 则
\[
  \ell'_{n_k}+\ell''_{n_k}\leqslant x_{n_k}+y_{n_k}.
\]
上式两边令 $k\to\infty$, 取极限即得
\[
  \underline{\lim_{n\to\infty}}\,x_n
  +
  \underline{\lim_{n\to\infty}}\,y_n
  =
  \lim_{k\to\infty}(\ell'_{n_k}+\ell''_{n_k})
  \leqslant
  \lim_{k\to\infty}(x_{n_k}+y_{n_k})
  =
  \underline{\lim_{n\to\infty}}\,(x_n+y_n).
\]
取 $\{x_n\}$ 的子列 $\{x_{n_k}\}$, 使得 $\lim\limits_{k\to\infty}x_{n_k}=\underline{\lim\limits_{n\to\infty}}\,x_n$. 由定理 2.5.2(2) 知
\[
  \underline{\lim_{n\to\infty}}\,(x_n+y_n)
  \leqslant
  \underline{\lim_{k\to\infty}}\,(x_{n_k}+y_{n_k}).
\]
再取 $\{y_{n_k}\}$ 的收敛子列 $\{y_{n'_k}\}$. 又由定理 2.5.2(2) 及 (3) 得
\[
\begin{aligned}
  \underline{\lim_{k\to\infty}}\,(x_{n_k}+y_{n_k})
  &\leqslant
  \underline{\lim_{k\to\infty}}\,(x_{n'_k}+y_{n'_k})\\
  &=
  \lim_{k\to\infty}x_{n'_k}+\lim_{k\to\infty}y_{n'_k}\\
  &=
  \underline{\lim_{n\to\infty}}\,x_n
  +
  \overline{\lim_{n\to\infty}}\,y_n.
\end{aligned}
\]

作为本节的结束, 我们给出两个例子来说明如何用上、下极限的理论来研究序列的敛散性。

\noindent\textbf{例 2.5.3}\quad 给定序列 $\{x_n\}$ 及常数 $\alpha\geqslant2$, 令
\[
  y_n=x_n+\alpha x_{n+1},\qquad n=1,2,\cdots,
\]
证明: 序列 $\{y_n\}$ 收敛的充分必要条件是序列 $\{x_n\}$ 收敛。

\noindent\textbf{证明}\quad 充分性是显然的, 下证必要性。

首先, $\{y_n\}$ 收敛蕴涵着它有界, 从而存在 $M>0$, 使得 $|y_n|\leqslant M$ 和 $|x_1|\leqslant M$. 现假定 $|x_n|\leqslant M$, 则有
\[
  |x_{n+1}|\leqslant\frac1\alpha(|y_n|+|x_n|)\leqslant\frac{2M}{\alpha}\leqslant M.
\]
由归纳法原理知, 对一切的正整数 $n$, 有 $|x_n|\leqslant M$, 即 $\{x_n\}$ 有界. 记 $\ell=\underline{\lim\limits_{n\to\infty}}\,x_n,\ h=\overline{\lim\limits_{n\to\infty}}\,x_n$, 则 $\ell$ 和 $h$ 均为有限数。

令 $\lim\limits_{n\to\infty}y_n=b$, 在等式 $x_n=y_n-\alpha x_{n+1}$ 两边分别取上、下极限, 并且利用本章习题第 32 题的结论和定理 2.5.3 的 (2), 可得
\[
  \begin{cases}
    h=b-\alpha\ell,\\
    \ell=b-\alpha h.
  \end{cases}
\]
由此可得 $\ell+\alpha h=h+\alpha\ell$, 从而有 $\ell=h$. 这样, 由定理 2.5.2 的 (3) 知, 序列 $\{x_n\}$ 收敛。

\noindent\textbf{例 2.5.4}\quad 证明序列 $\{\sin n\}$ 发散。

\noindent\textbf{证明}\quad 对 $\forall k\in\mathbb{N}$, 令
\[
  n_k=\left[2k\pi+\frac{3\pi}{4}\right],
  \qquad
  m_k=\left[(2k+1)\pi+\frac{3\pi}{4}\right],
\]
则有
\[
  2k\pi+\frac\pi4<n_k<2k\pi+\frac{3\pi}{4},
  \qquad
  (2k+1)\pi+\frac\pi4<m_k<(2k+1)\pi+\frac{3\pi}{4},
\]
从而
\[
  \sin n_k>\frac{\sqrt2}{2},
  \qquad
  \sin m_k<-\frac{\sqrt2}{2}.
\]
这样, 我们有
\[
  \underline{\lim_{n\to\infty}}\,\sin n
  \leqslant
  \underline{\lim_{k\to\infty}}\,\sin m_k
  \leqslant
  -\frac{\sqrt2}{2}
  <
  \frac{\sqrt2}{2}
  \leqslant
  \overline{\lim_{k\to\infty}}\,\sin n_k
  \leqslant
  \overline{\lim_{n\to\infty}}\,\sin n.
\]
因此, 序列 $\{\sin n\}$ 发散。

\section*{习题二}

1. 试将有理数集 $\mathbb{Q}$ 排成一个序列。

2. 试问以下陈述是否可作为序列极限的定义, 并说明理由:

(1) 对 $\varepsilon=\dfrac1{10^{100}}$, $\exists N>0$, 当 $n>N$ 时, 有 $|x_n-a|<\varepsilon$;

(2) 对 $\forall\varepsilon_k=\dfrac1k(k\in\mathbb{N})$, $\exists N>0$, 当 $n>N$ 时, 有 $|x_n-a|<\varepsilon_k$;

(3) $\exists M>0,\ \forall\varepsilon>0,\ \exists N>0$, 当 $n>N$ 时, 有 $|x_n-a|<M\varepsilon$。

3. 用序列极限的定义证明:

(1) $\lim\limits_{n\to\infty}\dfrac{\cos n}{n}=0$;\qquad
(2) $\lim\limits_{n\to\infty}\dfrac{n}{n^3+1}=0$;

\end{document}
